\documentclass[parskip=full]{article}

%%%%%%%%%%%%%%%%%%%%%%%%%%%%%%%%%%%%%%%%%%%%%%%%%
% Basic includes
%%%%%%%%%%%%%%%%%%%%%%%%%%%%%%%%%%%%%%%%%%%%%%%%%

\usepackage{arydshln}
\usepackage{array}
\usepackage{makecell}
\usepackage{tikz}
\usepackage{amsmath}
\usepackage{amsfonts}
\usepackage{verbatim}
\usepackage{mathrsfs}
\usepackage{amssymb}
\usepackage{amsthm}
\usepackage{bbm}
\usepackage{bm}
\usepackage[capposition=top]{floatrow}
\usepackage[dvipsnames]{xcolor}
\usepackage{ulem}
\usepackage[labelfont=bf]{caption}
\usepackage{subcaption}
\usepackage{graphicx}
\usepackage{float}
\usepackage{natbib}
\usepackage{booktabs}
\usepackage{longtable}
\usepackage{multirow}
\usepackage{rotating}
\usepackage{pdflscape}
\usepackage[title]{appendix}
\usepackage{footmisc}
\usepackage{bigdelim}
\usepackage{tabularx}
\usepackage{xspace}
\newcommand{\A}{\ensuremath{\mathcal{A}}\xspace}
\newcommand{\B}{\ensuremath{\mathcal{B}}\xspace}
\newcommand\pa[1]{\ensuremath{\left(#1\right)}}
\usepackage{epigraph}
\usepackage{setspace}
%\onehalfspacing
\setstretch{1.5}

%%%%%%%%%%%%%%%%%%%%%%%%%%%%%%%%%%%%%%%%%%%%%%%%%%%
% Microtype
%%%%%%%%%%%%%%%%%%%%%%%%%%%%%%%%%%%%%%%%%%%%%%%%%%%
\usepackage[activate={true,nocompatibility},final,tracking=true,kerning=true,spacing=true,factor=1100,stretch=10,shrink=10]{microtype}
% activate={true,nocompatibility} - activate protrusion and expansion
% final - enable microtype; use ``draft'' to disable
%tracking=true, kerning=true, spacing=true - activate these techniques
% factor=1100 - add 10% to the protrusion amount (default is 1000)
% stretch=10, shrink=10 - reduce stretchability/shrinkability (default is 20/20)

%%%%%%%%%%%%%%%%%%%%%%%%%%%%%%%%%%%%%%%%%%%%%%%%%%%
% Hyperref setup
%%%%%%%%%%%%%%%%%%%%%%%%%%%%%%%%%%%%%%%%%%%%%%%%%%%
\usepackage{url}
\makeatletter
\g@addto@macro{\UrlBreaks}{\UrlOrds}
\makeatother
\usepackage{hyperref}
\hypersetup{
  colorlinks,%
  citecolor=black,%
  filecolor=black,%
  linkcolor=blue,%
  urlcolor=black
}

%%%%%%%%%%%%%%%%%%%%%%%%%%%%%%%%%%%%%%%%%%%%%%%%%%%
% Font setup
%%%%%%%%%%%%%%%%%%%%%%%%%%%%%%%%%%%%%%%%%%%%%%%%%%%
\usepackage{mathpazo} % math & rm
%\linespread{1}        % Palatino needs more leading (space between lines)
\usepackage{parskip}
\usepackage[scaled]{helvet} % ss
\usepackage{courier} % tt
\usepackage{natbib}

\usepackage{filecontents}

\normalfont
\usepackage[T1]{fontenc}

\usepackage{titlesec}
\titleformat{\section}
{\normalfont\fontsize{12}{17}\bfseries}
{\thesection}
{1em}
{}

\titleformat{\subsection}
{\normalfont\fontsize{11}{16}\bfseries}
{\thesubsection}
{1em}
{}

%%%%%%%%%%%%%%%%%%%%%%%%%%%%%%%%%%%%%%%%%%%%%%%%%
% macros
%%%%%%%%%%%%%%%%%%%%%%%%%%%%%%%%%%%%%%%%%%%%%%%%%
\newcolumntype{C}[1]{>{\centering\let\newline\\\arraybackslash\hspace{0pt}}m{#1}}
\newcommand{\indicator}[1]{\mathbbm{1}_{\left\{ {#1} \right\} }}
%\newcommand{\resultsloc}{../code/python/output/}
\newcommand{\subfigwidth}{0.475\textwidth}
\newcommand{\subfigwidthtwo}{0.375\textwidth}
\newcommand{\figwidth}{0.6\textwidth}
\DeclareMathOperator*{\argmax}{arg\,max}

\setlength\epigraphwidth{0.9\textwidth}
\setlength\epigraphrule{0pt}
\renewcommand{\epigraphflush}{center}
\usepackage{etoolbox}
\makeatletter
\patchcmd{\epigraph}{\@epitext{#1}}{\itshape\@epitext{#1}}{}{}
\makeatother

%%%%%%%%%%%%%%%%%%%%%%%%%%%%%%%%%%%%%%%%%%%%%%%%%
% Page and paragraph settings
%%%%%%%%%%%%%%%%%%%%%%%%%%%%%%%%%%%%%%%%%%%%%%%%%
\setlength{\parskip}{0.0pc} \setlength{\parindent}{0.25in}
\setlength{\topmargin}{-3pc}
\setlength{\textheight}{9in}
\setlength{\oddsidemargin}{0pc} \setlength{\evensidemargin}{0pc}
\setlength{\textwidth}{6.5in}
\makeatletter
\@fpsep\textheight
\makeatother

\newenvironment{fignote}{\begin{quote}\footnotesize}{\end{quote}}
\def\sym#1{\ifmmode^{#1}\else\(^{#1}\)\fi}
\renewcommand{\baselinestretch}{1.5}

% Formats the notes at the bottom of tables and figures.
\newenvironment{fignote2}[1]
{
\par
\medskip
\begin{minipage}{#1}
\footnotesize
}
{
\end{minipage}
}

% Formats the notes at the bottom of tables and figures.
\newenvironment{fignote3}[1]
{
%\par
\begin{minipage}{#1}
\footnotesize
}
{
\end{minipage}
}

%%%%%%%%%%%%%%%%%%%%%%%%%%%%%%%%%%%%%%%%%%%%%%%%%
% Title page
%%%%%%%%%%%%%%%%%%%%%%%%%%%%%%%%%%%%%%%%%%%%%%%%%
%%%%%%%%%%%%%%%%%%%%%%%%%%%%%%%%%%%%%%%%%%%%%%%%%
% Title page
%%%%%%%%%%%%%%%%%%%%%%%%%%%%%%%%%%%%%%%%%%%%%%%%%

\author{Shahar Rotberg\footnote{University of Toronto. Corresponding author: \url{shahar.rotberg@mail.utoronto.ca}.} \and Lin Zhang\footnote{CMHC. Contact: \url{lzhang@cmhc-schl.gc.ca}. The opinions expressed herein do not necessarily represent the views of CMHC.}}

\title{\vspace{-1.8cm}The Welfare Consequences of Rent Control\footnote{We thank the seminar participants at University of Wisconsin-Madison, University of Windsor, and Amy Sun, Joseph B. Steinberg, Yahong Zhang, Ruoyu (Steven) Chen, Chenggang (Mark) Zhou, and John Baker for helpful and insightful comments. We are grateful to Joseph B. Steinberg for kindly allowing us to use his high-powered computer for this paper.}}

\pagestyle{plain}

\begin{document}

\maketitle

\setstretch{1.5}
\begin{abstract}
\onehalfspacing
\noindent We build a general equilibrium housing model in which renters sort into controlled and uncontrolled rental markets operating simultaneously, or become homelessness. Rent control insures against rent fluctuations and is redistributionary. However, to benefit from it, renters must initially pay a premium and their units deteriorate due to lack of maintenance. We calibrate our model to Toronto data, and show that although homelessness declines as rent control expands, it is optimal to have both markets. The uncontrolled market lets certain renters avoid the premium. The controlled market lets certain other renters avoid risk and benefit from redistribution and cheap housing.
\\[1.5ex]
% words=100

\noindent \textbf{JEL Classification:} E17; R31; R52. 

\noindent \textbf{Keywords:} General equilibrium, housing markets, rent control, homelessness.
\end{abstract}
\setstretch{1.5}

\section{Introduction}

Although tenancy decontrol is a common type of rent control in many cities around the world, the economic literature has only studied tenancy control (i.e., price ceilings). Price ceilings place restrictions on rent increases landlords can impose on a sitting tenant as well as upon unit turnover to a new tenant (e.g., New York City, Vancouver, Montréal, and Paris). Tenancy decontrol, on the other hand, places restrictions on rent increases on a sitting tenant but allows landlords to set rents freely when tenants change (e.g., Toronto, San Francisco, Los Angeles, Washington D.C., Singapore, and Hong Kong). 

In this paper, we develop a general equilibrium housing model that nests both tenancy decontrol and price ceilings. Our baseline is a model with tenancy decontrol, which we calibrate to data from Toronto, where tenancy decontrol is the rent control policy in place. We then ask if switching from our calibrated baseline to rent control that is based on price ceilings would be welfare-enhancing. Furthermore, we study how equilibrium prices and welfare would change in our baseline model with tenancy decontrol, following fully decontrolling and fully controlling the rental market. Lastly, we search for the welfare-maximizing degree of rent control in our baseline model, and compare it to using price ceilings with the same degree of rent control. 

What are the key differences between price ceilings and tenancy decontrol? The first difference relates to their ability to allocate controlled units to households with the highest willingness to pay. Price ceilings allow new tenants to benefit from immediate discounted rents relative to free market rents. However, this creates a shortage of rental housing because the quantity demanded exceeds the quantity supplied at the set price ceiling. This means that rental units are allocated through lotteries, and thus do not necessarily go to the tenants with the highest willingness to pay. On the other hand, because tenancy decontrol allows landlords to reset rents when tenants changes, market prices allocate controlled units to the renters with the highest willingness to pay. However, the price mechanism means that new tenants in controlled units must pay a premium relative to free market rents in order to secure their unit. This, in turn, means that it takes time until a new tenant renting a controlled unit begins to benefit from rent control.

The second difference relates to each rent control policy's ability to allocate controlled units to households with the highest marginal utility. Because housing cannot be consumed in arbitrarily small amounts, some low-income households, who have the highest marginal utility, may be shut out the rental market. These households may become homeless or resort to living with their parents. Price ceilings provide immediate rent discounts, and thus potentially remove barriers to entry into the controlled rental market for some low-income households with high marginal utilities. On the other hand, with tenancy decontrol renters must pay a premium in order to secure a controlled unit, which allocate controlled units to households with the highest willingness to pay. However, this may shut some households with the high marginal utilities out of the controlled market.  

We build a model to explore these tensions. 

This one would be good at the end of the intro: We unify tenancy control and decontrol in a general equilibrium model and compare the equilibrium effects of these two. We set the model of tenancy decontrol as our baseline model and compare its equilibrium with that of the model of tenancy control with lotteries. We also studied the equilibrium effects of the baseline model with full rent control and decontrol. In this way, we provided a unified model to study all types of existing rent controls. 



What market failures does rent control help address? 1. There is no insurance market to protect households from rent volatility. 2. Housing cannot be consumed at infinitesimal small amounts; thus, some households are shut out of the housing market, while other spend more on housing than they would if they could perfectly adjust their housing consumption, thus leading to suboptimal outcomes. Rent control attempts to address these failures by stabilizing rents and reducing housing costs to enable more households to spend more optimally on housing.  

Cities usually have no control over taxes and SAH. They can control local rents.

Many cities around the world have rent control laws with their main purpose to keep housing costs down. Even though in most cities controlled and uncontrolled rental markets operate simultaneous, there is no theory that explains how rents in these markets are determined in equilibrium. In the absence of a theory, it is difficult to assess whether rent control laws actually achieve their goal, through which channels, and whether they improve welfare. On the one hand, rent control erodes renters' housing costs during their tenancy and protects them against large fluctuations in their rents from one year to the next. On the other hand, rent control distorts rental prices by affecting the intensive and extensive margins of rental demand and by reducing rental maintenance. We provide a theory to explain how equilibrium rents are determined when controlled and uncontrolled markets operate simultaneously, and use our theory to assess the welfare implications of each of the aforementioned channels in general equilibrium.

To understand how rents are determined in equilibrium when uncontrolled (i.e., free market) and controlled markets operate in parallel, we develop a theory in which renters sort into rental markets or homelessness and then embed our theory into a general equilibrium model. Our first departure from the literature of heterogeneous-agent general equilibrium housing models is that we allow renters to choose whether to rent in the controlled market, the free market, or homelessness. In the controlled market, units are rent-controlled during a renter's tenancy; rents in those units decline in real terms the longer a renter stays in the same unit, but these units also deteriorate due to lack of maintenance by landlords, which is our second departure from the general equilibrium literature. When a tenant leaves a rent-controlled unit, landlords can command new market rents. Units in the free market are uncontrolled; landlords keep rental units in this market well-maintained, but rents in those units do not decrease in real terms with tenancy duration and are subject to stochastic fluctuations, which is our third departure from the literature. Due to market failure, renters cannot insure themselves against the financial risk they are exposed to in the free market. Renters sorting into homelessness do not pay rent, but receive little housing services, which we think of as shelters local governments provide. 

Our set-up has three main equilibrium consequences. The first is that in order to benefit from real decreases in rents over time and the insurance rent control provides against rent fluctuations, renters must initially pay a premium to rent in the controlled market. This is consistent with \citet{early-2000} and \citet{cjq-2023} who find that in New York City new tenants renting controlled units are more likely to actually have negative rental discounts. The second consequence is that as the share of the controlled market declines, competition for rent-controlled units drives up the premium. The third consequence is that high-income households who can bear the financial risk associated with the free rental market and who expect to become homeowners quickly choose to rent in the free market to avoid paying the premium. In contrast, low-income households whose risk tolerance is low and who expect to remain in the rental market for a long time choose to pay the premium to reap the benefits of rent control. This is consistent with the findings in \citet{gl-2003} for New York City, as well as our findings for Toronto, that renters who live in uncontrolled units earn, on average, income that is roughly 20\% higher than renters who live in controlled units. This is also consistent with our findings that renters in the controlled market are more than twice as likely as renters in the free market to stay more than 5 years in the same rental, as we show in Table \ref{tab:tenure_and_length_of_stay}.\footnote{In their Table A.7, \citet{cjq-2023} also show that in New York City, renters in the controlled market stay twice long in the same unit as renters in the uncontrolled market.}

Our theory and model set-up highlight the tradeoffs we discussed above, and thus allows us to examine the welfare effect of each. Rent control in our framework provides two benefits. In the presence of a market failure to insure against rent fluctuations, the first benefit is the provision of insurance against rent volatility.\footnote{Using data from the Rental Market Survey (RMS), we find that in Toronto, free market units' rents are 3.4 times more volatile than controlled units' rents.} The second benefit is redistribution from landlords to renters; rent control lowers tenants' rents over time.\footnote{\citet{zhang-gou} find that in Toronto, \textsl{ceteris paribus}, rent control reduced rents of controlled units by 11.3\% relative to otherwise-equivalent uncontrolled units. \citet{cjq-2023} produce similar findings for New York City.} On the other hand, rent control creates two distortions. First, rent control distorts rental prices by affecting the extensive and intensive margins of rental demand. Second, rent control reduces landlords' upkeep of their units, and thus controlled units deteriorate over time.\footnote{\citet{murray_et_al-1991} find that in Los Angeles, rent control laws similar to those in Toronto, reduced landlords' incentives to maintain their controlled units, causing them to deteriorate at an annual rate of 0.55\%.} We assess the magnitude of each these channels and its overall affect on newborns' welfare and homelessness.

The rest of our model has rather standard features. Households choose whether to rent or own a dwelling, or become homeless if they cannot live with their parents. Homeowners have access to long-term mortgage contracts. Working-age households earn labor income by supplying labor inelastically and receive a pension after retiring. Households' financial assets yield capital income. Labor and capital incomes are taxed progressively. A competitive rental company with convex marginal costs supplies rental housing. As in \citet{fallis-1985}, we model the rental housing supply after Toronto's rental market structure, where old construction is under rent control, but new construction is not. In this set up, the free market rent is the marginal revenue of the rental company, and thus determines the optimal quantity of rental housing the rental company supplies. See Section \ref{sec:rental_company} for more details. A competitive housing construction sector with convex marginal costs supplies new housing, which are added to the depreciated housing stock. 

The model also includes federal and local governments. The federal collects consumption and income taxes and uses the proceeds to pay pensions to retired households. The municipal government collects property taxes and uses the revenues to provide a public good that does not enter household utility. Because, in reality, rent control is under the control of local governments, we require budget balance only at the local level. We achieve budget balance by adjusting the property tax rate, which is the tax instrument under the control of local governments.

We calibrate our model to data from Toronto, Canada, where in 2016 72\% of rental units were controlled and the rest were uncontrolled. Most of our calibration is standard; we pick parameter values to match statistics like the homeownership rate, percentage of renters' income spent on rent, and so on. There are several targets, however, that are more unique to our calibration. The first two targets relate to the positive impacts of rent control. The first determines how much financial risk there is in the free market; we use data from the Rental Market Survey (RMS) to set rent fluctuations in the free market. The second governs how much redistribution rent control achieves; we set the annual rate at which rents in the controlled market decrease in real terms for tenants who remain in the same unit to 1.37\%, using data from the 2016-2021 Rental Market Survey (RMS) the 2016 Canadian Census. 

On the other hand, we use three other statistics to determine the negative impacts of rent control. The first statistic controls annual rate of deterioration of controlled units, which we set to 0.55\% based on estimates from \citet{murray_et_al-1991}. For our second statistic, we follow \citet{rs-2023} and set the rental supply price elasticity to 0.3 to match a renters' property tax incidence in big cities, such as New York City, Boston, Paris, and so on. The rental supply elasticity dictates the degree of price distortions rent control causes because it affects how much rents change in equilibrium when rental demand shifts. The lower the elasticity, the larger the distortions. Lastly, as in \citet{rs-2023}, we pick the smallest rental unit to match the percentage of renters who spend over 50\% of their income on rent. This statistics determines how much the price distortions rent control creates affect the welfare of low-income renters. The greater the impact on low-income renters, the greater the impact on aggregate welfare.   

We use the calibrated model with two rental markets as the benchmark to compare with our experiments, and compare long-run stationary equilibria with the benchmark. We adopt the concepts in \citet{ralws-1971} and define welfare as the ex-ante expected utility of newborns, which is a common practice in the literature \citep{conesa-et-al,gkkoc}

Our first experiment compares the benchmark with the long-run equilibrium in which rental markets are fully decontrolled. Full decontrol generates a tradeoff between initial rents of new tenants and average rents. On the one hand, decontrolling would reduce the value of staying longer in the same rental unit and encourage homeownership. Thus, decontrolling would raise house prices and reduce total rental demand. The lower rental demand would lower free market rents and initial rents paid by new tenants. On the other hand, the average rent would rise because the decline in free market rents would be insufficient to offset the substantially lower average rents of benchmark-controlled units. Despite the higher average rents, welfare would rise and all newborns would benefit. In particular, low-income newborns would benefit more from lower initial rents than they would lose from higher average rents. High-income newborns would also benefit, but for a different reason; they expect to become homeowners quickly, and thus tend to rent in the free market, which would become cheaper following full decontrol. Figure \ref{fig:diagram_full_decontrol} provides further visual intuition.

Our second experiment compares the benchmark with the long-run equilibrium in which all rental units are under control. With full rent control, the average rent in the controlled market would become the marginal rent of the rental company. Since the average controlled rent is 10.2\% lower than the free market rent in the benchmark, the initial rent new tenants pay under full rent control would have to increase significantly for the rental company to supply the benchmark-quantity of rental housing. As this would unfold to bring the rental market back to equilibrium, renters' lifecycle rent profile would become expensive. In response, more renters would switch to homeownership, and those that remain renters would reduce their rental demand. The lower total rental demand would dominate the higher homeownership, reducing total housing demand and house prices. Overall, the vast majority of newborns would be worse off and welfare would decline. This is because the higher rents new tenants would face would outweigh the benefits of slightly lower house prices. Figure \ref{fig:diagram_full_control} provides further visual intuition.

In our third experiment, we search for the long-run welfare-maximizing degree of rent control. We change the share of rental housing mandated under rent control and compute welfare in the long-run equilibrium under each case. We find that no degree of rent control is optimal, implying full rent decontrol is the optimum. Why? When less of the rental market is controlled, there are fewer renters who optimally remain in the same unit and more renters who switch to owning. As homeownership rises, total demand for rental housing declines, which in turn lowers equilibrium free market rents. As a result, initial rents paid by new tenants decline, and because this benefits newborns, welfare rises. 

Our last experiment computes the transition from the benchmark to the long-run equilibrium under full decontrol. We model the transition following an unexpected permanent shock. Upon shock, renters living in large rental units previously protected by rent control, would no longer find it optimal to remain in those now-expensive units. This would reduce rental demand and rapidly reduce rents. This, in turn, would lower initial rents paid by newborns, and thus immediately raise welfare and throughout the transition. To the best of our knowledge, this is the first general equilibrium analysis of the transition from an economy with rent control to an economy without rent control.

Beyond the welfare analyses we expounded upon above, this paper contributes empirically and theoretically to the research on rent control. Empirically, we document new stylized facts about rental markets; the income gap between renters that rent in the uncontrolled market and renters renting in the controlled, and the rent premium that new tenants must pay to access controlled units. Theoretically, we build a model in which renters choose which rental market to rent in, and show that our calibrated model generates equilibrium outcomes similar to the ones we observe in the data, thus offering a new theory of how controlled and uncontrolled rental markets operate concurrently in equilibrium.

%The rest of the paper proceeds as follows. We discuss the contribution of this paper to the literature in Section \ref{sec:related_lit}. Next, we introduce the institutional background on rent control laws in Toronto in Section \ref{sec:evidence} and  document the stylized facts using several datasets in Section \ref{sec:evidence}. Then, we present our model in Section \ref{sec:model} and calibrate the model in Section \ref{sec:calibration}. Finally, we conduct experiments in Section \ref{section:experiments} and conclude in Section \ref{sec:conclusion}.

\section{Related literature}\label{sec:related_lit}

\textcolor{red}{SR: need to add discsussions of the following papers: Glaeser and Luttmer (2003) - comparing income of renters in two markets, Autor et al (2014) - end of rent control in MA, Basu and Emerson (2000, 2003) - good articles to use as reason for our model to be reasonable for many cities in the US (and consider their modelling contributions), Monras and Montalvo (2023) - empirical on rent control in Catalonia. Frankena (1975), Raymon, (1983), Sims (2007), Downs (1988); Moon and Stotsky (1993) - last 5 about quality affect. I think most of these papers are empirical and do not assess the welfare effects of each channel that we model. Plus we highlight premium to rent in controlled units, and two channels that are totally ignored in the literature: price distortions in free market and insurance. We also quantify the impact of each of these channels.}

Several strands of the rent control literature are related to our analysis. Prominent economists such as \citet{hayek-1929} and \citet{fs-1946} studied how rent freezes affect the rental housing supply and the deadweight loss arising from rent controls. \citet{fallis-1985} use a static demand and supply model with controlled and free markets interacting with each other. These analyses use static models and do not include an ownership market, and thus are unsuitable for exploring the interaction between rental markets and house prices, or the dynamic lifecycle aspects of rent control.

\citet{fmn-2022} use a heterogeneous-agent general equilibrium model to study the impact of housing policies such as rent vouchers on housing markets. However, their rent-controlled market resembles social affordable housing rather than the rent control we study here. Assignment of social affordable housing units usually has a significant luck element and confers immediate rent discounts to new tenants as modelled in \citet{fmn-2022}. In our model, renters choose which market to rent in, implying that to clear markets, renters choosing to rent in the controlled market must initially pay a premium. \textcolor{red}{discuss the tradeoffs they focus on regarding rent control and compare to ours.}

\textcolor{red}{can probably remove this part}
We support our theory with data for Toronto Census Metropolitan Area (CMA), and embed our theory in a general equilibrium model to explain how rents in different markets are endogenously determined in the presence of rent control. Our findings are consistent with \citet{early-2000} and \citet{cjq-2023} who estimate that in New York City tenants in the controlled market are systematically more likely to pay rent premiums during the first few years of tenancy.\footnote{See Section 3 in \citet{early-2000} and Figure B4 in \citet{cjq-2023}.} Both \citet{early-2000} and \citet{cjq-2023} were puzzled by their findings of rent premiums paid by renters in the controlled market and \citet{early-2000} even mechanically set them to zero. Our theory and model offer a solution to this puzzle by showing that rent premiums in the controlled market are an equilibrium outcome of renters choosing which rental market to rent in. 

\citet{sims-2007} empirically evaluates the effects of rent control on rents and the quantity of rental housing supplied using the sudden stop of rent control in Massachusetts, but to the best of our knowledge, our work is the first to provide a complete theoretical analysis of the transition out of rent control. In contrast to \citet{sims-2007}, \citet{diamond-2019} evaluate the welfare impact of expansions in rent controls in San Francisco using a dynamic, neighborhood choice model. Their model features many desirable features, such as tenant buyouts, the potential for biased renters' expectations about the future, and endogenous landlord responses. However, it does not feature general equilibrium elements that we focus on, such as endogenous rents, feedback effects between free market rents and controlled rents, or interactions between the rental and homeownership markets. As \citet{diamond-2019}, we also find rent controls expansions would raise average rents and hurt newborns.

\textcolor{red}{add discussion about papers we were referred to by referees}

\section{Institutional background on rent control in Toronto, Canada}\label{sec:insitutional_background}
The National Housing Act introduced rent controls in Canada in 1944 during World War II. This “first-generation” of rent control froze nominal rents and was inspired by the rent freeze in European countries during World War I intended to protect veterans from evictions and rent increases. The rent freeze seemed justified during the war, but many economists believed that maintaining it after had significant negative consequences, including reduced rental supply, labor mobility, and property maintenance \citep{arnott-1995}. As a result, the Ontario provincial government repealed the “first-generation” of rent controls in 1949.  

In 1975, during the oil crisis, the Ontario government re-introduced rent controls as part of the Residential Premises Rent Review Act. These new controls, classified as “second-generation” rent controls, differed in two ways from “first-generation” rent controls. First, they allowed mild nominal rent increases once a year based on mandated guidelines. Moreover, courts could also grant increases beyond the statutory guidelines in cases of major increases in the cost of supplying rental housing, financial losses of landlords, or required capital expenditures. Second, they exempted luxury accommodation, government owned buildings, social housing, and private rental units built after January 1976.

Since 1975, the Ontario rent control system has been evolving with two major changes. The first is intermittent exemptions. In 1985, the Residential Rent Regulation Act removed exemptions for previously exempted units and newly built units. Under this new change, both new buildings and buildings constructed after 1976 were subject to rent controls. In 1992, rent control with exemptions for new construction was again restored by The Rent Control Act, which exempted all rental buildings built after November 1, 1991. However, in April 2017, the government of Ontario removed exemptions by extending rent controls to rental units built after November 1, 1991. In 2018, the Ontario government enacted new legislation exempting only buildings built, or first occupied, after November 15, 2018.

The second major change in the Ontario rent control system came in 1998, and was marked by tenancy deregulation as part of the Ontario Tenant Protection Act. This Act allowed landlords of rent-controlled units to increase rents by more than prescribed by the rent increase guidelines, but only when signing a new contract with a new tenant. Otherwise, legislation compelled them to follow the rent increase guidelines. This type of rent control is sometimes referred to as tenancy rent control, partial rent control, or “third-generation” rent control \citep{arnott-2003}. The Ontario government has not made any changes since. 

\textcolor{red}{review of other cities with rent control}

\section{Model}\label{sec:model}

The model features overlapping generations of finitely-lived households, competitive rental companies, competitive housing construction companies, and governments. Households in each age cohort are heterogeneous in terms of labor market productivity, housing size, mortgage debt, and financial assets, and housing status (own, rent, live with their parents, or homelessness). Renters also differ in the rental market they are renting (i.e., controlled or free market). Renters in the controlled market also differ in the number of years they have lived in the same rental unit. Renters in the free market are subject to rent shocks. 

Before retirement, households work in the labor market and after retirement they receive pensions from the government. Pensions depend on households' average lifetime labor income. Renters choose which market to rent in and how long to remain in their units. Homeowners pay proportional property taxes and maintenance costs each period, and mortgages are modelled as long-term contracts, as in \citet{rs-2023}. As in \citet{rs-2023}, before the age $J_{p}$, households can decide to live with their parents rent-free at the expensive of receiving less housing services than they would in the rental market. As in \citet{az-2024}, households whose are older than age $J_{p}$ cannot live with their parents, and thus if they cannot or decide not the rent in the rental market, they become homeless, which means that they receive little housing services for free. 

Competitive rental companies with convex marginal costs operate simultaneously in the controlled and free rental markets. Competitive construction companies with increasing marginal costs invest in the production of new housing, which is added to the depreciated housing stock. 

We include federal and local governments. The federal government provides pensions to retired households through taxes on income and consumption. The provincial and municipal governments provide services to residents by imposing taxes on property. Because rent control laws are the mandate of local governments, we only require that the local government balance its budget in equilibrium. 

Because the economy is a city, the interest rate, $r$, is exogenous \citep{sommeretal-2013,arslanetal-2015,sommersullivan-2017,kaplanetal-2019}.\footnote{The assumption of a small open economy is prevalent in the quantitative housing literature. See, for example, \citet{sommeretal-2013}, \citet{arslanetal-2015}, \citet{floettotetal-2016}, \cite{sommersullivan-2017}, \citet{kaplanetal-2019}, and \citet{kkk-2021}. We allow households to save in bonds and also purchase stock in the rental and construction companies, but a no-arbitrage condition requires households to be indifferent between these two options. Because we do not require the asset market to clear, we do not need to explicitly keep track of the dividends associated with these stock holdings. Another interpretation is that the rental and construction companies are owned by foreigners, as in \citet{floettotetal-2016}.}  

\subsection{Demographics}\label{subsection:demographics}

Households are born at age $j=1$, retire from the labor force at age $J_{R}$, and live $J$ periods. They do not face mortality risk before age $J$, and die with certainty after age $J$. Households have the following preferences:

\begin{equation*}
    u_{j}(c,h(o,f,\iota))=\xi_{j}^{\sigma}\frac{\left(c^{1-\gamma} \left(h(o,f,\iota) \right)^{\gamma}\right)^{1-\sigma}}{1-\sigma},
\end{equation*}

\noindent where $c$ represents non-housing consumption, $h$ is housing services consumption, and $f=1$ denotes households renting in the free market, $f=0$ denotes households renting in the controlled market.\footnote{For owners, households living with their parents or homeless households, the value of $f$ does not matter, and so we set it to $1$.} $\iota$ denotes the number of years a renter in the controlled market has lived in the same unit, $o=1$ denotes homeowners and $o=0$ denotes renters, households living with their parents, and households who are homeless. $\gamma$ is the utility weight on the consumption of shelter services, $\sigma$ is the household's constant relative risk aversion, and $\xi_{j}$ is an equivalence-scaling to capture variation in household size over its lifecycle.\footnote{See \citet{lm-1980}, \citet{wz-2010}, \citet{xu-2013}, and \citet{rs-2023} for references on the equivalent scale, $\xi_{j}$.} We define $h(o,f,\iota)$ in Section \ref{subsection:housing_choice} below. Households discount the future at a rate of $\beta$. On non-housing consumption, households pay a proportional tax, $\tau_{c}$.

As in \citet{cagetti-denardi-2006}, households who die receive a warm-glow from bequests they leave behind, captured by $\omega_{B}(q)=\mu \frac{q^{1-\sigma}}{1-\sigma}$. $q$ stands for the household's net worth at the end of their life, which we define below, and $\mu$ affects the strength of the bequest motive.

\subsection{Heterogeneity}

Households are heterogeneous in: financial assets, $a\in \mathcal{A}=\mathbb{R}_+$; whether they live with their parents, $v\in \mathcal{V}=\{0,1\}$; whether they are homeless, $b \in \mathcal{B}=\{0,1\}$; their tenure (own or rent), $o \in \mathcal{O}=\{0,1\}$, in case they get their housing services through the housing market; the rental market they rent in if they receive their rental services from the rental market, $f \in \mathcal{F}=\{0,1\}$; years they lived in the same unit if they rent in the controlled market, $\iota \in \mathcal{I}=\mathbb{N}$; shock to free market rent if they rent in the free market, $z \in \mathbb{Z}$; house size, $h\in \mathcal{H}=\mathbb{R}_{++}$; mortgage debt, $m \in \mathcal{M}=\mathbb{R}_+$; idiosyncratic labor fixed effect $\chi \in \mathcal{X}$, and stochastic labor market productivity shock, $e \in \mathcal{E} = \mathbb{R}_{++}$. The measure of age-$j$ households over endowments is $\Psi_{j}(a,v,b,o,f,\iota,z,h,m,\chi,e)$. For brevity, we mostly use $s$ to denote a household's state $(a,v,b,o,f,\iota,z,h,m,\chi,e)$.

\subsection{Income and income taxes}\label{subsection:income}

Working-age households ($j<J_{R}$) supply labor inelastically. Age-$j$ households with labor market productivities, $\chi$ and $e$, earn labor income $y_{\ell,j}(\chi,e)=\chi e\zeta_j$. The fixed-effect $\chi$ is constant throughout a household's life and is drawn from an unconditional distribution, $\hat{F}(\chi)$. $e$ is an idiosyncratic shock that follows a Markov process that is also conditional on whether the household is homeless or not, $F(e,e'|b)$, where $b=1$ means that the household is homeless and $b=0$ means that they are not homeless. In particular, as in \citet{az-2024}, households that are homeless are more likely than households that are not homeless to get stuck in a low stochastic productivity shock. $\zeta_j$ is a common life-cycle profile. Retired households receive a pension, $y_{R}(\chi,e)$, which depends on labor productivity at retirement. Retirees with a high level of productivity are more likely to have had higher productivities throughout their lives, and thus receive larger pensions. 

Households can also generate income from their savings, $y_{k}(a)=ar$, where $r$ is an exogenous interest rate. We define taxable income as:

\begin{equation*}\label{eq:taxable_labor_income}
  y^{x}_{j}(\chi,e,a) =y_{\ell,j}(\chi,e)+\nu y_{k}(a),
\end{equation*}

\noindent where $\nu$ is the taxable portion of capital income. We apply progressive taxes to taxable income as follows:

\begin{equation}\label{eq:income_tax}
  y^{\tau}_{j}(\chi,e,a) = y^{x}_{j}(\chi,e,a)-\tau_{\ell}y^{x}_{j}(\chi,e,a)^{\psi},
\end{equation}

\noindent where $\tau_{\ell}$ determines the average tax and $\psi$ affects the degree of progressivity \citep{hsv-2020}.

As in \citet{kkk-2021}, each period, living households receive endowments of financial wealth equal to the bequests left by households who have died in the previous period. Surviving households of all ages receive financial endowments equal to $\vartheta (\chi,e)$, which captures the correlation between bequests and income.

\subsection{Housing}\label{subsection:housing_choice}

Households can choose to rent or own a dwelling, $o \in \mathcal{O}=\{0,1\}$, live with their parents $v \in \mathcal{O}=\{0,1\}$ if $j<J_{p}$, or become homeless $b \in \mathcal{O}=\{0,1\}$ if $j \geq J_{p}$. Households who choose to rent or own can also choose their dwelling size, $h$. The smallest rental unit is smaller than the smallest housing unit for purchase, defined as $0<\underline{h}_{r}<\underline{h}_{o}$ \citep{gervais,sommeretal-2013,sommersullivan-2017,kaplanetal-2019}. Households who choose to live with their parents receive housing services equal to $\underline{h}_{p}$, and households choose to become homeless get housing services equal to $\underline{h}_{h}$, where $0<\underline{h}_{h}<\underline{h}_{p}<\underline{h}_{r}$. Households who choose to live with their parents, ($v=1$), or who choose to be homeless ($b=1$) receive their housing services for free. However, the city incurs a cost $\theta_{b}$ to support each homeless household. A household deciding to own a house of size $h$, pays the purchase price $p h$. Buying and selling houses entail one-time proportional transaction costs, $\tau_{b} p h$ and $\tau_{s} p h$, respectively. Each period, all housing units are taxed at a rate $\tau_{p}$ and depreciate in value at a rate $\delta_h$. 

Renters decide whether to rent in the controlled market or the free market, where $f \in \{0,1\}$; $f=0$ stands for renting in the controlled market and $f=1$ represents renting in the free market. Renters who choose to rent housing services $h$ in the free market rent pay $p_{r,1}h$. Their rent is independent of their tenancy duration, and it is subject to uncertain shocks, $z$, which follow a distribution $G(z)$. With probability $G(z=x)$, renters who spend $p_{r,1}h$ in the free market today, wake up in the next period with additional rent payments equal to $x p_{r,1}h$. $x$ is an independent shock that can take on either negative and positive value. The average per-unit rent in the free market is thus $\int_\mathcal{Z}(1+z)p_{r,1}\mathrm{d}G(z)$, which we denote by $\overline{p}_{r,1}$. 

On the other hand, rent in the controlled market is certain and decreases with tenancy duration. Similarly to \citet{fmn-2022}, the longer the tenancy duration in the controlled market, the cheaper the real rent becomes. As in \citet{fmn-2022}, controlled rent is a function of tenancy duration as follows:

\begin{equation}\label{eq:rent_in_controlled_market}
\frac{p_{r,0}}{(1+\kappa_{1})^{\iota}} h,
\end{equation} 

\noindent where $h$ is the condition of the controlled unit at the time the new tenant moved in. $p_{r,0}$ is the rental price paid by new tenants, and $\kappa_{1}$ is the rate at which rent control reduces real rents for each additional period tenants remains in the same unit. $\iota$ is the number of years they have lived in that unit. Whereas landlords of free market units fully offset the depreciation of their units each period, landlords of controlled units do not maintain these units as well. Thus, controlled units deteriorate over time, reducing the amount of housing services tenants derive from those units \citep{gl-1990,la-2018}. The amount of housing services tenants derive in the controlled market is thus as follows:

\begin{equation}
    h(o,f,\iota)=\frac{1}{(1+(1-o)(1-f)\kappa_{2})^{\iota}} h.
\end{equation}
 
\noindent $\kappa_{2}$ is the rate at which controlled units deteriorate for each additional period a household remains renting the same unit. 

Figure \ref{fig:diagram_rent_path} qualitatively illustrates rents paid by renters as a function of the market in which they rent and their tenancy duration.\footnote{Because a rental unit's deterioration is lower than the decrease in its real rent, we abstract from rental unit deterioration to make our point more crisp.} The downward sloping blue line shows rents paid by renters living in controlled units. When they enter a new rental unit, their rent is $p_{r,0}$, but as their tenancy duration increases, their rent decreases as in equation (\ref{eq:rent_in_controlled_market}). On the other hand, the horizontal red line shows average rents paid by renters living in the free market, $\overline{p}_{r,1}$. The average rent remains constant regardless of their tenancy duration, but their actual rent can fluctuate around the average rent. 

Figure \ref{fig:diagram_rent_path} suggests that, in equilibrium, the average rent paid in the free market, $\overline{p}_{r,1}$, must be lower than the rent paid by renters in the controlled market early on in their tenancy, $p_{r,0}$. However, the average rent in the controlled market $\overline{p}_{r,0}$ is lower than $\overline{p}_{r,1}$. What explains this equilibrium outcome? If initial rents paid by new tenants in the controlled market were lower than the average rents paid by renters in the free market, $\overline{p}_{r,1}$, all renters would choose to rent in the controlled market. To clear the markets, renters in the controlled market must initially pay a premium. Renters anticipating to remain only a few periods in the same unit, and who have enough risk tolerance for rent fluctuations, would choose to rent in the free market to avoid paying the premium. In contrast, renters expecting to remain many periods in the same unit, and whose risk tolerance for rent volatility is low, would pay the premium and enjoy protection from rent risk and lower long-term rents.

\subsection{Mortgages}\label{subsection:morgtages}

Besides rent control, mortgage contracts that households can obtain is another quantitatively important economic factor affecting households' decisions to buy or rent. We follow \citet{rs-2023} and model mortgages, $m$, as long-term debt with origination costs, prepayment penalties, and borrowing restrictions. We set the mortgage interest rate to $r_{m}=r+\Delta_m$, where $r$ is the return on financial assets, and $\Delta_m$ is the mortgage debt interest rate premium \citep{sommeretal-2013,sommersullivan-2017,kaplanetal-2019}. Households cannot deduct mortgage interest payments from their taxable income because mortgage interest payments are not tax-deductible in Canada. To originate a new mortgage, households must pay a fixed cost, $\omega_{1}$. At the time of new mortgage origination only, the loan to value ratio (LTV) must not exceed a fraction $\lambda_{m,1}(ph)<1$, which depends on the value of the house purchased, $ph$:

\begin{equation}\label{eq:LTV}
\frac{m}{p h} \leq \lambda_{m,1}(ph).
\end{equation}

\noindent A constraint of payment to income ratio (PTI) also limits mortgage debt at origination: 

\begin{equation}\label{eq:PTI}
\frac{\tau_{p}ph+(r_{m}+\upsilon_{1})m}{y_{\ell,j}(\chi,e)+y_{k}(a)} \leq \lambda_{m,2}.
\end{equation}

\noindent $\upsilon_{1}$ is the fraction of debt paid every year on average over the life of a mortgage, and $\lambda_{m,2}$ is the maximum PTI ratio allowed. Constraint (\ref{eq:PTI}) says that property taxes, interest payments, and the average fraction of the principal paid cannot exceed the fraction, $\lambda_{m,2}$, of their pre-tax labor and capital incomes. 

Each period, a mortgage holder must pay at least a fraction $\upsilon_{1}$ of their total mortgage debt, and are allowed to pay up to a fraction $\upsilon_{2}>\upsilon_{1}$ of the debt without a penalty. Any amount paid above $\upsilon_{2}m$ is subject to a prepayment penalty as follows: 

\begin{equation}\label{eq:prepayment_penalty}
\tau_{pp}(m,m')=\omega_{2}\mbox{max}\left\{(1-\upsilon_{2})m-m',0\right\}.
\end{equation}

\noindent If a mortgage holder pays less than $\upsilon_{1}m$, we treat them as having originated a new mortgage, and thus they must incur the origination cost, $\omega_{1}$, as well as satisfy constraints (\ref{eq:LTV})--(\ref{eq:PTI}) again.

\subsection{Dynamic program}\label{subsection:household_optimization}
Given its state variables, a household chooses how much to save ($a'$), whether to own or rent ($o' \in \left\{0,1\right\}$), whether to live with their parents ($v' \in \left\{0,1\right\}$), or whether to become homeless ($b \in \left\{0,1\right\}$). Owners and renters choose how much housing to consume ($h'$). Renters also choose whether to rent in the controlled or free market ($f' \in \left\{0,1\right\}$, where we set $f'=1$ if $o'=1$, $v'=1$, or $b'=1$), and renters renting in the controlled market also decide whether to remain in their rental unit one period longer ($\iota'$). Homeowners decide how much mortgage debt to take on ($m'$) and whether to originate a new mortgage ($n$). The dynamic program representing this problem is:

\begin{align}
  \label{eq:vf}
  V_{j}(s)=\max_{\{c,a',v',b',o',f',\iota',h',n,,m',q'\}}\ \biggl\{u_{j}\left(c,h'\right(o',f',\iota'\left)\right)+\indicator{j=J}\omega_{B}(q')+
  \indicator{j<J}\beta \int_{\mathcal{E}\times \mathcal{Z}}V_{j+1}(s')dF(e,e'|b')dG(z)\biggr\}.
\end{align}
Subject to
\begin{flalign}
  & c(1+\tau_{c})+a'+or_{m}m+(1-b')(1-v')(1-o')\left[(1-f')\frac{p_{r,0}}{(1+\kappa_{1})^{\iota'-1}}h'(0)+f'p_{r,1}h'\right]+\label{eq:bc} \\ & \quad (1-b)(1-v)(1-o)fzp_{r,1}h+o\left(\tau_{p}+\delta_{h}+\indicator{o'=0 \vee h'\neq h}\tau_{s}\right)ph + o'(1+\tau_{b}\indicator{o=0 \vee h'\neq h})ph'+ \notag \\ & \quad\omega_{1}n+(1-n)\tau_{pp}(m,m') \indicator{o'=1 \land j=J}ph(\tau_{p}-\delta_{h}-\tau_{s})= \notag \\ & \quad y_{\ell,j}(\chi,e)+y_{k}(a)-y^{\tau}_{j}(\chi,e,a)+a+o(ph-m)+o'm'+\vartheta(\chi,e) \notag \\
  & a' \geq 0 \label{eq:financial_wealth} \\
  & h' \geq (1-b')(1-v') o' \underline{h}_{o}+(1-b')(1-v')(1-o')\underline{h}_{r}+b'\underline{h}_{h}+v'\underline{h}_{p} \label{eq:min_housing}\\
  & v'=0 \mbox{ if } j\geq J_{p} \label{eq:can_live_with_parents}\\
  & b'=0 \mbox{ if } j<J_{p} \label{eq:can_be_homeless}\\
  & h'(o',f',\iota')=\frac{1}{(1+(1-o')(1-f')\kappa_{2})^{\iota'-1}}h' \label{eq:housing servies_controlled}\\
  & n \in \left\{0,o'\indicator{j<J_{R}}\right\} \label{eq:origination} \\
  & nm' \leq \lambda_{m,1}(ph') p h' \label{eq:LTV_1} \\
  & n (\tau_{p}ph'+(r_{m}+\upsilon_{1})m')\leq \lambda_{m,2}(y_{\ell,j}(\chi,e)+y_{k}(a))\label{eq:PTI_1} \\
  & (1-n)m'\leq (1-\upsilon_{1})m \label{eq:settle_debt} \\
  & q' = a' + o' \left[p h'-m'\right] \label{eq:networth}
\end{flalign}

\noindent Household utility includes three components: a flow utility from housing and non-housing consumption; a warm-glow flow from bequests left behind by households who die; and future value is discounted by $\beta$.\footnote{We include labor market productivity shocks for retirees as well as working-age households in our formulation of the dynamic program for notational brevity. Writing the dynamic program this way allows us to use a single equation to represent both retired and working-age households' problems.} Rentals units rented by households in the controlled market deteriorate over time due to lack of maintenance by landlords, and this deterioration shows up as a reduction in the consumption of housing services.

Equation (\ref{eq:bc}) is the household's budget constraint. The left-hand side contains expenditures: consumption and consumption taxes, $c(1-\tau_{c})$; saving, $a'$; interest payments for mortgage holders, $r_{m}m$; the cost of renting a unit in the controlled market for renters renting in that market, $\frac{p_{r,0}}{(1+\kappa_{1})^{\iota'-1}}h'$, where $h'$ is the condition in which the controlled unit was when the renter moved into it; the cost of renting in the free market for renters renting in that market, $p_{r,1}h'$; and rent shocks, $z p_{r,0}h$, for households who rented in the free market in the previous period. The left-hand side also includes homeownership-related costs: property taxes, depreciation costs, and selling costs for homeowners switching their house, $(\tau_{p}+\delta_{h}+\indicator{o'=0\vee h'\neq h}\tau_{s})ph$; home purchasing costs for homeowners, $(1+\tau_{b}\indicator{o=0\vee h'\neq h})ph'$; origination costs for households originating a new mortgage, $\omega_{1}$; prepayment penalty costs for mortgage holders paying more than $\upsilon_{2}$ of their mortgage, $\tau_{pp}(m,m')$; and lastly, property taxes, depreciation and selling costs for households who reach age $J$ and choose to end their lives as homeowners. The right-hand side of the budget constraint contains resources available to spend: after-tax labor and capital incomes, $y_{\ell,j}(\chi,e)+y_{k}(a)-y^{\tau}_{j}(\chi,e,a)$; financial assets, $a$; housing equity for owners, $(ph-m)$, mortgage funds raised, $m'$, for mortgage holders; and bequests $\vartheta(\chi,e)$ that depend on the household's labor productivity.

Constraint (\ref{eq:financial_wealth}) says that households do not have access to unsecured debt. Constraint (\ref{eq:min_housing}) requires all owners consume at least $\underline{h}_{o}$ units of housing, all renters consume at least $\underline{h}_{r}$, homeless households to consume $\underline{h}_{h}$, and households living with their parents consume $\underline{h}_{p}$. Constraint (\ref{eq:can_live_with_parents}) says that households whose age is equal to greater than $J_{p}$ cannot live with their parents. Constraint (\ref{eq:can_be_homeless}) says that households whose age is strictly less than $J_{p}$, can live with their parents and thus cannot be homeless. Constraint (\ref{eq:housing servies_controlled} says that renters in the controlled market consume less housing services as they stay longer in the same unit. Constraint (\ref{eq:origination}) says that a household can only originate a new mortgage if it is of working-age. Constraints (\ref{eq:LTV_1}) and (\ref{eq:PTI_1}) are the LTV and PTI constraints at mortgage origination, respectively. Constraint (\ref{eq:settle_debt}) says that mortgage holders must repay at least a fraction $\upsilon_{1}$ of their mortgage principal. Lastly, equation (\ref{eq:networth}) says that the amount of bequests left by households who die, are equal to the sum of their financial and net housing wealth.

We denote the policy functions for consumption, savings, living with parents, becoming homeless, tenure, rental market choice, years living in the same unit, housing, mortgage debt, mortgage origination, and bequests, by $c_{j}(s)$, $a'_{j}(s)$, $b'_{j}(s)$, $v'_{j}(s)$, $o'_{j}(s)$, $f'_{j}(s)$, $\iota'_{j}(s)$, $h'_{j}(s)$, $m'_{j}(s)$, $n_{j}(s)$, and $q'_j(s)$, respectively.

\subsection{Rental company}\label{sec:rental_company}

We follow the rental housing supply theory in \citet{fallis-1985} and embed it in \citet{garrigaeetal-2009}'s rental housing supply structure. Let $i=0$ denote the controlled market and $i=1$ the free market. For each rental housing unit in market $i$, $x_{i}=1$, the rental company obtains average rent, $\overline{p}_{r,i}$, pays property taxes, $\tau_{p}$, and incurs depreciation costs, $\delta_{h}$. 

To supply a total of $X'=\sum^{1}_{i=0} X'_{i}$ units of rental housing, the rental company also incurs management costs equal to $\frac{\theta_{1}}{\theta_{2}}X'^{\theta_{2}}$, where $\theta_{1}$ determines the average cost to manage the stock of rental housing $X'$, and $\theta_{2}\geq 1$ determines price elasticity of the rental housing supply. 

The rental company must supply $\mbox{min}\{\underline{X}_{0},X'\}$ in the controlled market, with $\underline{X}_{0} \geq 0$. This means that if the total quantity of rental housing supplied is $X'\leq \underline{X}_{0}$, the market is completely under control. In this case, the rental company's marginal revenue is the average rent in the controlled market, $\overline{p}_{r,0}$. However, if $X'>\underline{X}_{0}$, the rental company's marginal revenue is the average free market rent, $\overline{p}_{r,1}$. We also refer to the rental company's marginal revenue as the marginal rent. In this case, $\underline{X}_{0}$ denotes the quantity of controlled units supplied and $X'-\underline{X}_{0}$ the quantity supplied in the free market.

The rental company buys houses, without friction, at the house price, $p$, and each period decides how much rental housing to supply in the controlled market, $X'_{0}$, and how much to supply in the free market $X'_{1}$. The rental company has the following problem:

\begin{align}\label{eq:rental_company_problem}
V(X)=\underset{X'_{0},X'_{1}} {\text{$\mbox{max} $}}\left\{\sum_{i=0}^{1} \overline{p}_{r,i}X'_{i}-\frac{\theta_{1}}{\theta_{2}}X'^{\theta_{2}}-p X'+p(1-\delta_{h}-\tau_{p})X+\frac{1}{1+r}V(X')\right\}.
\end{align}
Subject to
\begin{align}
& X'_{0}=\mbox{min}\{\underline{X}_{0},X'\} \label{eq:supply_restriction} \\
& X'_{1} \geq 0 \label{eq:supply_positive} \\
& X'=\sum^{1}_{i=0} X'_{i}. \label{eq:total_supply} 
\end{align}

\noindent Equation \ref{eq:supply_restriction} says that the quantity of rental housing supplied in the controlled market is the smaller element between $\underline{X}_{0}$ and $X'$. Equation \ref{eq:supply_positive} says that the quantity supplied in the free market cannot be negative. Equation \ref{eq:total_supply} says the total quantity supplied is the sum of the quantities supplied in both markets. Whenever $X'>\underline{X}_{0}$ the supply function is:

\begin{align}\label{eq:rental_supply}
X'^{\star}=\left[\frac{1}{\theta_{1}}\left(\overline{p}_{r,1}-p\left(\frac{r+\delta_{h}+\tau_{p}}{1+r}\right)\right)\right]^{\frac{1}{\theta_{2}-1}}.
\end{align}

\noindent In equilibrium, when $X'^{\star}\geq \underline{X}_{0}$, then $X'^{\star}_{0}=\underline{X}_{0}$ and $X'^{\star}_{1}=X'^{\star}-\underline{X}_{0}$. If $X'^{\star}<\underline{X}_{0}$ or if rent control laws apply to all rental units, then $\overline{p}_{r,0}$ replaces $\overline{p}_{r,1}$ in equation \ref{eq:rental_supply}, $X'^{\star}=X_{0}$, and $X'^{\star}_{1}=0$. 

\noindent The partial derivatives of $X'^{\star}$ with respect to $p$ and $\overline{p}_{r,1}$, respectively, show the optimal quantity supplied:

\begin{enumerate}
    \item Decreases with the housing price and increases with rent if $\theta_{2}>1$, which is the case in our calibration.
    \item Is independent of the housing price and rent if $\theta_{2}=1$; in this case, rental supply is perfectly elastic.
\end{enumerate}

One could argue that as time goes on, the share of the controlled market decreases because as rent controlled units depreciate they may be replaced with new construction not subject to rent control. Thus, how can the share of the rent controlled market be preserved in stationary equilibrium? In order to maintain a constant controlled market share, governments can increase the year of construction beyond which units are exempted from rent control. This is exactly what the Ontario government has done over the years. A similar policy is also in place in Quebec, where new construction is exempted from rent control for only five years.

\subsection{Housing construction company}\label{sec:construction_company}

A competitive construction company builds new housing \citep{rotberg-2022}. For its production of new housing, $I$, the construction company receives revenues $pI$, but incurs costs captured by the following technology:

\begin{equation*}
TC(I)=\frac{\epsilon_{1}}{\epsilon_{2}}I^{\epsilon_{2}}\label{eq:total_cost_structure},
\end{equation*}

\noindent where $\epsilon_{1}$ determines the average cost of housing investment and $\epsilon_{2}$ controls the price elasticity of housing investment. The construction company chooses investment to maximize profits:

\begin{equation*}
\label{eq:construction_company_profits}
\underset{I} {\text{$\mbox{max}$}}{\left\{p I-TC(I)\right\}}.
\end{equation*}

\noindent The solution produces the following relationship between housing investment and the price of housing:

\begin{equation}
\label{eq:construction_company_first_order_condition}
I^{\star}=\epsilon_{1} p^{\epsilon_{2}-1}.
\end{equation}

\subsection{Aggregation}\label{subsection:aggregation}

As we mentioned in the general description of the model section, we model federal and local governments. The federal government collects income and consumption taxes and provide pensions to households. The provincial and municipal governments collects property taxes, and provides services to residents and shelter services and supports to the homeless. Since rent control is policy determined at local levels, we only require budget balance at that lower level of government:
\begin{align}\label{eq:gbc}
  G + \sum_{j=1}^J \int_\mathcal{S} \theta_{b} b'_{j}(s)\mathrm{d}\Psi(s)  = \tau_{p} p H.
\end{align}

\noindent Government expenditures on services to residents and homeless support must equal to revenues from property taxes. When calibrating the model to the current tax code in Toronto, we choose the consumption level of public goods, $G$, to ensure a balanced budget. In our experiments, we keep $G$ constant, and adjust $\tau_{p}$ to balance the budget. 

$S=A\times V\times B\times O \times F \times I \times Z \times H\times M\times E$ denotes a typical subset of the state space $\mathcal{S}$; and $A$, $V$, $B$, $O$, $F$, $I$, $Z$, $H$, $M$, $X$ $E$ denote typical subsets of $\mathcal{A}$, $\mathcal{V}$, $\mathcal{B}$, $\mathcal{O}$, $\mathcal{F}$, $\mathcal{I}$, $\mathcal{Z}$, $\mathcal{H}$, $\mathcal{M}$, $\mathcal{X}$, and $\mathcal{E}$, respectively.

Aggregate market housing demand must equal aggregate housing supply:

\begin{equation}
\label{eq:housing_clearing_condition}
  \sum_{j=1}^J \int_\mathcal{S} (1-b'_{j}(s))(1-v'_{j}(s)) h'_{j}(s) \mathrm{d}\Psi(s) = H',
\end{equation}

\noindent and the housing stock law of motion is 

\begin{equation}\label{eq:housing_stock_lom}
H'=I+(1-\delta_h)H.
\end{equation}


\noindent Equation (\ref{eq:housing_stock_lom}) is a standard law of motion for housing, which says that the next period stock of housing is equal to investment in housing plus the depreciated stock of housing. In each market, aggregate rental housing demand must equal aggregate rental housing supply:

\begin{equation}
\label{eq:rental_housing_clearing_condition}
  \sum_{j=1}^J \int_\mathcal{S} (1-b'_{j}(s))(1-v'_{j}(s)) (1-o'_{j}(s))\indicator{f'_{j}(s)=i}h_{j}(s) \mathrm{d}\Psi(s) = X'_{i} \quad i \in \{0,1\}.
\end{equation}

The distributions of households evolve according to the laws of motion,
\begin{align}\label{eq:lom2}
\Psi_{1}(\mathrm{S})=\int_{\mathcal{E} \times \matchcal{X}} \indicator{(0,0,0,0,0,0,0,\underline{h}_{r},0,X,E)\in \mathrm{S}}\mathrm{d} \hat{F}(X)\tilde{F}(E),
\end{align}

\begin{align}
\Psi_{j+1}(S) &=\int_{\mathcal{S}} \hat{F}(X)\hat{F}(X)F(E,e|b'_{j}(s))G(z)M_{j}(s)\mathrm{d}\Psi_{j}(s),\ j=1,\ldots,J-1\label{eq:lom1}
\end{align}

\noindent where,

\begin{align}
    M_{j}(s) = \indicator{[a'_{j}(s)\in A] \land [v'_{j}(s)\in V] \land [b'_{j}(s)\in B] \land [o'_{j}(s)\in O] \land [f'_{j}(s)\in F] \land [\iota'_{j}(s)\in I] \land [h'_{j}(s) \in H] \land [m'_{j}(s) \in M]}.\notag
\end{align}

Lastly, total bequests left by decedents must equal total wealth endowments:
\begin{equation}
\label{lom:agg_bequests}
      \int_{\mathcal{S}} q'_{J}(s) \mathrm{d}\Psi_{J}(s) =  \sum_{j=1}^J \int_{\mathcal{S}} \vartheta (\chi,e)\ \mathrm{d}\Psi_j(s).
\end{equation}
Like \citet{kkk-2021}, we pick $\vartheta$ such that this condition holds in all of our experiments. This means that changes in $\vartheta$ between equilibria depend on the proportional changes in aggregate bequests.

\subsection{Equilibrium}
An equilibrium in the model is a sequence of aggregate prices and quantities, $\left\{H',I,\{X'_{i},\overline{p}_{r,i}\}^{1}_{i=0},p,q\right\}$, a sequence of value and policy functions, $\biggl\{V_{j}(\cdot),c_{j}(\cdot),a'_{j}(\cdot),v'_{j}(\cdot),b'_{j}(\cdot),o'_{j}(\cdot),f'_{j}(\cdot)\iota'_{j}(\cdot),h'_{j}(\cdot),m'_{j}(\cdot),n_{j}(\cdot),q'_{j}(\cdot)\biggl\}_{j=1}^J$, and a sequence of measures, $\left\{\Psi_{j}(\cdot)\right\}_{j=1}^J$, that:
\begin{enumerate}
\item Solve the household's dynamic problem (\ref{eq:vf}) subject to constraints (\ref{eq:bc})--(\ref{eq:networth});
\item Solve the rental company's problem (\ref{eq:rental_company_problem}) subject to (\ref{eq:supply_restriction})--(\ref{eq:total_supply});
\item Satisfy the construction company's first-order condition, (\ref{eq:construction_company_first_order_condition});
\item Satisfy the government's budget constraint, (\ref{eq:gbc});
\item Clear the housing market,
(\ref{eq:housing_clearing_condition}) subject to the housing stock law of motion (\ref{eq:housing_stock_lom});
\item Clear both rental housing markets,
(\ref{eq:rental_housing_clearing_condition});
\item Satisfy the laws of motion for the distribution of households, (\ref{eq:lom1})--(\ref{eq:lom2});
\item And satisfy the law of motion for bequests, (\ref{lom:agg_bequests}).
\end{enumerate}

\section{Calibration and model performance}\label{sec:calibration}

Calibration follows two steps. First, we calibrate certain parameters using data for Toronto and estimates from the literature. Second, we calibrate the values of the remaining parameters so that the model generates statistics to match Toronto data. Table \ref{tab:calibration_external} lists the externally-set parameter values and sources. Table \ref{tab:calibration_internal} lists the internally-set parameter values and their targets. In Section \ref{appendix:data} of the appendix, we expand on how we estimated certain statistics using various data sets. Third, in Section \ref{subsection:non_targeted_statistics}, we show how the model performs on several dimensions we did not target. 

\subsection{Assigned parameter values}

\subsubsection{Demographics, preferences, and income}
A model period corresponds to 1 year. Households live from age 26 to age 85, and retire at age 66, implying retirement age $J_{R}=41$ and lifespan $J=60$. As in \citet{rs-2023}, we set the maximum age households can live with their parents, $J_{p}$, is set to 9, so that households aged 35 and above cannot live with their parents. We set households' relative risk aversion $\sigma=2$. The equivalent household size scales for each age group, $\{\xi_{j}\}_{j=1}^{J}$, are set based on 2016 Census data on the average household size for each age cohort. We scale household sizes by their square roots, as it is standard in the literature. The discount factor, $\beta$, the utility weight on housing services, $\gamma$, and the utility advantage to homeownership, $\eta$, are determined in the second stage of our calibration.

We take the labor productivity process from \citet{gkkoc}. We set the standard deviation of the fixed effect to $0.128$. The stochastic labor productivity shocks follow an AR(1) process in logs,

\begin{equation*}
  \log e' = \rho_{e} \log e + \epsilon_{e},\quad \epsilon_{e}\overset{\mathrm{i.i.d.}} \sim N(0,\sigma_{e}).
\end{equation*}

\noindent We set the persistence, $\rho_{e}=0.9$, as in \citet{gkkoc}. We set the standard deviation of labor productivity, $\sigma_{e}$, in the second stage. Newborns draw their labor productivity from an unconditional probability distribution associated with the ergodic distribution generated endogenously by our productivity processes. As in \citet{az-2024}, for households who choose to be homeless, the we mechanically set the probability of remaining in their stochastic labor productivity state, $e$, to be 20\% higher than what is implied by our stochastic labor productivity process above.

Using 2016 Census data, we set the deterministic lifecycle component of labor market productivity to: 
\begin{equation*}
\zeta_j = 0.923+0.09\times j-1.61\times 10^{-3}\times j^{2}, \mbox{ for } j\leq J_{R}
\end{equation*}

\noindent This formula comes from regressing average household income for each age from 26 through 65 on age and age$^{2}$. We set pension benefits, $\{y_{R}(e)\}_{e\in \mathcal{E}}$, using the replacement rate of gross individual earnings \citep{OECD-pension}. 

Surviving households' endowments are proportional to their stochastic labor income: $\vartheta(e):=\vartheta \chi e$. We choose $\vartheta$, which controls surviving households' wealth endowments, to satisfy equation (\ref{lom:agg_bequests}).

\subsubsection{Taxes, interest rates, and housing}

The taxable portion of capital income, $\nu$, is set to 50\% as in Canada's tax code. Using 2016 Census household-level data for Toronto, we estimate $\psi=0.9275$ in equation (\ref{eq:income_tax}). We set the consumption tax rate, $\tau_{c}$, to 13\% to equal the Harmonized Sales Tax (HST) rate in Ontario. We set the interest rate, $r$, to 2.77\% as estimated for Canada by \citet{bruneau-etal-2018}. 

We set our rent shocks $\left\{z\right\}_{\mathbb{Z}}=\left[-0.1227,-0.0607,0,0.0624,0.1284,0.2589\right]$ and the corresponding probabilities $G(z)=\left[0.16,0.12,0.44,0.08,0.16,0.04\right]$ using 1992–2016 data from the RMS on average rents in the controlled and uncontrolled market. Our data suggest that there is a higher risk of large rent increases than benefits of substantial rent decreases. Using Census and estimate from \citet{zhang-gou}, we set the rate at which real rents decrease with tenancy duration in the controlled market, $\kappa_{1}$, to equal to 1.365\%. We use \citet{murray_et_al-1991}'s estimate of the annual deterioration rate of controlled units in Los Angeles to set $\kappa_{2}=0.55\%$. The per-household cost to the city to support the homeless, $\theta_{b}$, is set to 82.3\% of the average household income. We set this parameter by dividing the City of Toronto total spending on homelessness supports by the of number of homeless individuals in the City of Toronto in 2021.\footnote{See \url{https://www.homelesshub.ca/community-profile/toronto} and \url{https://www.toronto.ca/legdocs/mmis/2021/ex/bgrd/backgroundfile-163804.pdf}.}

We set the property tax rate, $\tau_p$, to 0.44\% using information reported by Altus Group and 1992-2016 Property Sale and Assessment Data (PSAD). We set the housing depreciation rate to 0.6\% as estimated by \citet{rosenthal-2014}. Based on estimates by \citet{gm-2003}, we set buying costs, $\tau_{b}$, to 2\% of the purchase price, and selling costs, $\tau_{s}$, to 6\% of the sale price. We choose $\epsilon_{2}=1.89$ so that the price elasticity of housing supply is equal to $0.89$, as estimated by \citet{piaxao-2021} for Toronto. The level of the construction company's cost function, $\epsilon_1$, is a free parameter that we back out from the market clearing condition \eqref{eq:housing_clearing_condition}. Lastly, we set the weighted average rent in both rental markets, $\overline{p}_{r}$, to 5.24\% of the equilibrium house price, which we computed using 2016 data from the Canadian Real Estate Association (CREA).

\subsubsection{Mortgages}

Following banking regulations in Canada, we set the LTV constraint, $\lambda_{m,1}(ph)$, as follows:

\begin{equation*}
\lambda_{m,1}(ph) = \begin{cases}
0.8 &\text{if $ph\geq \$1M$}\\
0.9 &\text{if $\$500K\leq ph< \$1M$}\\
0.95 &\text{if $ph < \$500K$}.
\end{cases}
\end{equation*}

\noindent Mortgage holders with a model-equivalent house value above \$1M can borrow only up to 80\% of its value; mortgage holders with a model-equivalent house value between \$500,000--\$1M can borrow up to 90\% of the house value; and mortgage holders with a house value below \$500,000 can get a loan of up to 95\% of the house value. We set $\lambda_{m,2}=0.39$ so that the maximum allowable PTI is 39\% as is standard in the mortgage industry in Canada.\footnote{See \url{https://www.cmhc-schl.gc.ca/consumers/home-buying/calculators/debt-service-calculator}.}

For mortgage contracts, we first set the fixed mortgage origination cost, $\omega_{1}$, to 0.68\% of the average household income in Toronto based on data from the 2016 Census and Canada Mortgage and Housing Corporation's (CMHC) website.\footnote{See \url{https://www.cmhc-schl.gc.ca/en/professionals/industry-innovation-and-leadership/industry-expertise/resources-for-mortgage-professionals/how-much-will-my-home-really-cost}. Mortgage origination costs include the cost of appraisal ($\approx\$250$) and lawyer fees ($\approx\$500$).} We set the prepayment penalty, $\omega_{2}$, to 3 months of interests, a common practice in mortgage contracts in Canada. We follow \citet{rs-2023} for setting $\upsilon_{1}$ and $\upsilon_{2}$. We set the average fraction of mortgages paid each year, $\upsilon_{1}$, which is used in inequalities (\ref{eq:PTI}) and (\ref{eq:settle_debt}), to 1.7\% to reflect a 30-year amortization period, which is common in the mortgage industry in Canada. We set the threshold for prepayment penalty, $\upsilon_{2}$, to $1.15 \times \upsilon_{1}$ plus 15\% of the total debt and interest payments, which adds up to 18.8\%.\footnote{Standard mortgage contracts in Canada allow prepayments of 15\% of the principal and regular payments without penalty.} Finally, we set the interest rate premium on mortgage debt, $\Delta_{m}$, to 0.19\%, which is the average spread between the 5-year variable rate on mortgages and the return on a 10-year government bond from 2002 to 2022.\footnote{The mortgage data we used is not publically available. You can find 10-year Government of Canada bond yield here: \url{https://www.bankofcanada.ca/rates/interest-rates/canadian-bonds/}}

\subsection{Internally calibrated parameters}

After assigning the parameter values listed above, we still need to calibrate another eleven parameters. We list them in Table \ref{tab:calibration_internal}. The variables are the discount factor, $\beta$; the utility weight on the consumption of housing services, $\gamma$; the bequest motive, $\mu$; the average level of income taxation, $\tau_{\ell}$; the variability in the stochastic component of labor income, $\sigma_{e}$; the level of income taxation, $\tau_{\ell}$; the convexity of managing rental housing, $\theta_{2}$; the marginal managerial cost to rent ratio, determined by $\theta_{1}$; the rent a tenant pays in the first period after entering a rental unit under rent control, $p_{r,0}$; the smallest rental unit, $\underline{h}_{r}$; the smallest housing unit for purchase, $\underline{h}_{o}$; the housing services a household who lives with their parents gets, $\underline{h}_{p}$; and the housing services a households who is homeless gets, $\underline{h}_{h}$.\footnote{The construction cost scalar, $\epsilon_{1}$, is a free parameter chosen to equate housing supply to housing demand in the benchmark equilibrium. In counterfactual simulations, we hold $\epsilon_{1}$ fixed at its benchmark level.}

We jointly calibrate all these parameters to match eleven moments in Toronto data: the aggregate mortgage debt to aggregate house value of homeowners, which is 21.5\% according to the 2019 wave of the Survey of Financial Security (SFS); the percentage of total renters' income spent on rent, which is 22.5\% according to the 2016 Census; the ratio of aggregate bequests to aggregate net wealth, which is 2.2\%, based on our estimates using the 2019 SFS; the average labor income tax rate, which is 19.6\%, and comes from the 2016 Census; the Gini coefficient of labor income , which comes from the 2016 Census; the property tax incidence of renters for cities similar to Toronto, which is about 12\%;\footnote{We take the average property tax incidence of renters estimated by \citet{cy-1994}, \citet{susin-2002}, \citet{fack-2006}, and \citet{sy-2020}. \citet{cy-1994}, \citet{fack-2006}, and \citet{sy-2020} directly estimate the property tax incidence. \citet{susin-2002}, on the other hand, measures the percentage of housing vouchers captured by landlords. Because he finds that housing vouchers are fully captured by landlords, this means that the rental housing supply curve is perfectly inelastic, as they also show in their appendix, and the property tax incidence of renters is 0\%. See Section 2.1 in \citet{rs-2023} for a detailed discussion on this estimation method.} the free market rent to price ratio, which is 5.82\%, based on our calculations using the 2016 Census and 2016 CREA data; the percentage of the renters in the the controlled rental market is 71\%, which we compute using the 2016 Census. This determines the premium paid by new tenants in the controlled market: as the number of renters in the free market increases and the number of renters in the controlled market declines, fewer rental housing is available in the controlled market and the higher the equilibrium rent premium must be to clear rental markets. We also target the percentage of renters spending over 30\% of their gross income on rent, which comes from the 2016 Census and is equal to 41\%; the homeownership rate, which is 66.5\% according to the 2016 Census; the percentage of 26-34 year old households who live with their parents, which we estimate is 22.3\% using the 2016 Census; and lastly we target the percentage of households who are homeless, which is 0.26\% and comes from the ``Homeless Hub''.\footnote{See \url{https://www.homelesshub.ca/community-profile/toronto}.}

Although all parameters are determined jointly, in practice each parameter affects one moment more substantially than others. The discount factor, $\beta$, mostly determines the amount of debt homeowners take on: the lower the discount factor, the greater the debt homeowners optimally take on. The utility weight on housing, $\gamma$, influences the percentage of their total income renters' on rent: the greater $\gamma$, the more housing services households want to consume. The strength of the bequest motive, $\mu$, affects the bequest to net wealth ratio: the larger the bequest motive, the more bequests are left relative to net wealth. The average level of income taxation, $\tau_{\ell}$, determines the average income tax rate: the lower $\tau_{\ell}$, the more tax households pay on their incomes. The variability in the stochastic component of labor income, $\sigma_{e}$, determines the degree of inequality in labor income: the higher $\sigma_{e}$, the higher the labor income Gini coefficient. Furthermore, the marginal cost of managing rental housing, $\theta_{2}$, determines the price elasticity of the rental housing supply: the higher $\theta_{2}$, the more price-inelastic the rental supply, and the lower the property tax burden of renters. This target generates in our calibrated model a price elasticity of rental housing supply equal to 0.3. $\theta_{1}$ influences the free market rent to price ratio: the higher the rent to price ratio, the higher $\theta_{1}$. The initial rent paid for a controlled unit, $p_{r,0}$, controls the percentage of renters renting in the free market: the higher the initial rent paid in the controlled market relative to the free market, the larger the percentage renting in the free market. The smallest rental unit, $\underline{h}_{r}$, governs the percentage of renters spending over 30\% of their gross income on rent: the larger the smallest rental, the larger the percentage of renters spending over 30\% of their income on housing. The smallest housing unit for purchase, $\underline{h}_{o}$, controls the homeownership rate: the larger the the smallest housing unit for purchase, the more fewer households optimally choose to own a home. The amount of housing services conferred to households living with their parents, $\underline{h}_{p}$, impacts the percentage of 26-34 year old households living with their parents: the greater the housing services from living at home, the greater the percentage of households who live with their parents. Lastly, housing services obtained by homeless households, $\underline{h}_{h}$, impacts the homelessness rate: the greater the housing services obtained by homeless households, the greater the homelessness rate. 

\subsection{Model performance}\label{subsection:non_targeted_statistics}

We evaluate our model performance against aggregate, lifecycle, and tenure distribution statistics. Table \ref{tab:aggregate_non_targreted_statistics} shows the performance of our model on aggregate statistics that are not targeted in the calibration. The model produces a somewhat lower average owned-house value to average household income ratio relative to the data, as well as a somewhat lower percentage of owners with positive mortgage debt. Our data do not allow us to estimate the ratio of the average value of a rental unit to the average value of an owner-occupied unit in Toronto. Therefore, we compare it to \citet{chatterjee-eyigungor-2015}'s estimate for this ratio for the United States, which is somewhat lower than our model's. Using 2016--2021 RMS data, we estimate the average tenure duration in the same rental unit is 8.2 years. Our model predicts a slightly higher average tenure duration. Although we do not have comparable estimates in the data, we include two other statistics. The premium for new tenants in the controlled market is about 8.6\% in our model, and the average free market rent is 0.7\% higher than the average rent in the controlled market. Importantly, in our model renters in the free market earn 18.8\% more than renters in the controlled market, which is almost exactly what we observe in the 2016 Census. Without a means-testing allocation, only a model in which renters sort into the controlled and free market can generate the latter result.

The wealth distribution in our model is less concentrated than in the data. This is because our model does not feature rate-of-return heterogeneity, which \citet{fagereng} argue is paramount to generating a high concentration of wealth at the top of the wealth distribution.\footnote{Our model also does not feature heterogeneous rate of time preferences \citep{ks-1998} or Michael Jordan income shocks \citet{mj-shocks}, which could also explain its inability to generate a high wealth concentration.}

Figure \ref{fig:life_cycle} shows lifecycle outputs of the model relative to the data. Panel (a) shows the homeownership rate over the lifecycle in our model is slightly lower than the data for young and old age groups, and somewhat higher than the data for middle-aged households group. Panel (b) shows young households in our model have higher LTV ratios compared to young households in the data, whereas the reverse is true for old households. 

In Figure \ref{fig:rent_to_income_cdf} we show the rent to income cumulative frequency distribution. Our model generates similar rent to income ratios as observed in the data, although it has slightly fewer households with high rent to income ratios and slightly more households with low rent to income ratios. 

Lastly, Table \ref{tab:tenure_by_income_quintile} shows households' tenure by their position in the income distribution. The second column from the left is the lowest income quintile and the first column from the right is the highest income quintile. We break down tenure into two categories: the percentages of renters in each quintile renting in the controlled, and the percentage of renters in each quintile renting in the free market. Panel (a) shows these statistics in the data, whereas panel (b) shows the calibrated model's performance on the same statistics. Although the calibrated model cannot fully replicate the data, it can generate important features of the data for the purposes of our research: households in the lower quintiles of the income distribution are more likely to rent in the controlled market than those in the higher income quintiles.

\begin{comment}
Conferences in order of time:
1. Midwest Economics Association (MEA) in March 2023: https://mea.sites.grinnell.edu/ 
2. Midwest Macro Group in May 2023: https://krannert.purdue.edu/faculty/mcrucini/midwest-macroeconomics-group.php
3. Society for Economic Dynamics (SED) in June 2023: https://www.economicdynamics.org/
4. Canadian Macro Study Group (CMGS) in November 2023: https://canadianmacro.ca/
5. American Economic Association (AEA) in January 2024:
https://www.aeaweb.org/conference/

Planned presentations in universities:
1. Madison Wisconsin - November 2023
2. Ryerson - May 2023.
3. May reach out to McMaster and York as well, I have a few friends there, and I presented there before. 
\end{comment}

\section{Experiments}\label{section:experiments}

\paragraph{Experiments' overview} Using our calibrated model, we conduct four main experiments. In our first experiment we use our model to measure effects of three features in the model on the premium and welfare in the benchmark. We start with the impact of risk, move on to the impact of redistribution, and finish with the impact of deterioration in the controlled market. In the second experiment, we compare the benchmark (i.e., rent control with exemptions or partial rent control) with the long-run equilibrium without any rent control, which we refer to as full decontrol. In the third experiment, we compare the benchmark with the long-run equilibrium in which all rental housing is controlled, which we refer to as full control. In the last experiment, we search for the degree of rent control that maximizes welfare in the long-run equilibrium. In all our experiments, we keep the tax code fixed and balance the government's budget by adjusting the property tax rate. 

\paragraph{Welfare measure} Our main measure of welfare is the expected ex-ante utility of newborns \citep{ralws-1971,conesa-et-al,gkkoc}. It is equivalent to giving all newborns a fixed proportional non-housing consumption transfer. Formally, the non-housing consumption-equivalent welfare effect due to a change in policy is given by

\begin{equation}
  CE_{t} = \left[\frac{\sum_{j=1}^{k}\int_{\mathcal{S}}V^{\dagger}_{j}(s)\mathrm{d}\Psi^{\dagger}_{j}(s)-\sum_{j=1}^{k}\int_{\mathcal{S}}\omega_{B}(q'^{\star}_{J}(s))\mathrm{d}\Psi^{\dagger}_{j}(s)}{\sum_{j=1}^{k}\int_{\mathcal{S}}V^{\star}_{j}(s)\mathrm{d}\Psi^{\star}_{j}(s)-\sum_{j=1}^{k}\int_{\mathcal{S}}\omega_{B}(q'^{\star}_{j}(s))\mathrm{d}\Psi^{\star}_{j}(s)}\right]^{\frac{1}{(1-\gamma)(1-\sigma)}}-1,
\end{equation}

\noindent where index daggers denote counterfactual equilibrium objects and stars denote objects in the benchmark. We also compute the approval of the new policy change by measuring the mass of households whose welfare rises under the new policy change:
\begin{equation}
A = \sum_{j=1}^{k}\int_{\mathcal{S}} \indicator{V^{\dagger}_{j}(s)>V^{\star}_{j}(s)}\ \mathrm{d}\Psi^{\star}_{j}(s).
\end{equation}

\noindent For our main welfare and approval measures, we set $k=1$, which measure the welfare and approval of newborns. We also present average welfare gains and approval rates across all households by setting $k=J$. Because we compute approval using the benchmark distribution, we can also compute the approval of households who are renters in the benchmark and households who are homeowners in the benchmark. 

\paragraph{Measures of rent} Our analysis features several measures of rent listed below: 

\begin{enumerate}
    \item \underline{Average free market rent}: the equilibrium average rent in the free market.
    \item \underline{Average controlled rent}: the equilibrium average rent in the controlled market.
    \item \underline{Rent of new tenants}: the equilibrium initial rent new tenants pay when moving into a new rental unit. In the controlled market, new tenants pay more than tenants who have been under rent control for a while. In the free market, new tenants expect to pay the same average rent as tenants living in the same rental unit for a while.
    \item \underline{Average rent}: the weighted average of equilibrium rents in the controlled and free rental markets.
\end{enumerate}

\subsection{Price ceilings}

\subsubsection{Aggregate results}

In our first experiment we study the impact that three key features in our model have on welfare and the rent premium in the controlled market in our benchmark. We begin by examining the effect of rent volatility risk in the free market, then explore the impact of redistribution created by rent control, and finish with the impact of deterioration of controlled units. 

For each case, we eliminate the effect from our benchmark and then solve the model in partial equilibrium with that feature omitted. Comparing each partial equilibrium to the benchmark gives us the welfare effect of each feature in the benchmark. Then, for each case, we solve the full general equilibrium with the specific feature omitted. In general equilibrium, removing each feature would affect the premium in the controlled market, and by comparing it to the benchmark premium, we obtain the impact of each of the features on the premium. We show our results in Table \ref{tab:premium_breakdown}.

\subsection{Benchmark vs. full decontrol}

\subsubsection{Aggregate results} The second column of Table \ref{tab:aggregate_effects} presents changes in aggregate variables in the decontrolled equilibrium relative to the benchmark. Full decontrol would reduce the value of staying longer in the same rental unit and cut tenure duration by 28\%. This would encourage renters to own, reduce total rental demand, and lower free market rents by 6.5\%. The increase in homeownership would outweigh the decline in total rental demand, and thus total housing demand would increase and with it house prices. The higher house prices would increase property tax revenues and enable the government to reduce property taxes slightly.

How is it possible that free market rents would decrease but average rents would increase? Free market rents in the benchmark are 10.2\% higher than the average rent in the controlled market, and the share of the controlled market in the benchmark is almost 90\%. Thus, a 6.5\% decline in free market rents would not be sufficient to lower average rents, and the average rent would rise despite the decline in the quantity of rental housing supplied. More specifically, decontrolling would increase average rents paid on benchmark-controlled units, but reduce average rents paid on benchmark-uncontrolled units. Overall, the former would outweigh the latter. Figure \ref{fig:diagram_full_decontrol} provides a qualitative diagram with detailed explanations to visually demonstrate these effects.

Decontrolling would benefit all newborns and raise their welfare. Because newborns have low incomes, they would benefit more from the lower initial rents under full decontrol than from the lower average rents later in life in the benchmark. However, decontrolling would reduce average welfare because higher average rents would hurt older renters who have lived in controlled units for a while. Perhaps surprisingly, decontrolling would benefit about 70\% of benchmark-homeowners but only 30\% of benchmark-renters, totaling to a slight majority of households with welfare gains. 

Why do many homeowners gain? While the higher house prices would increase homeowners' net wealth and benefit them, they would also raise homeowners' ongoing proportional maintenance costs and hurt them. Moreover, the lower rents in the free market would benefit homeowners by protecting them in situations in which they would need to switch to renting due to a negative income shock or retirement. Overall, the benefits would outweigh the losses for most homeowners.

Why would a majority of renters lose? This is because most renters have lived in controlled units for a while, and decontrolling would result in rent increases for them.  

\subsubsection{Distributional results} 

Full rent decontrol would increase all newborns' welfare, but how would these gains be distributed? To answer this question, we break down newborns' welfare by labor productivity. Table \ref{tab:distributional_impact} presents welfare changes by labor market productivity. $e_{1}$ represents the lowest labor productivity, and $e_{7}$ the highest. All newborns would gain from the rent effect of decontrolling. Since newborns with low productivity have low incomes and assets early in life, they would benefit from trading off lower rents early in life in exchange for higher average rents later in life. Newborns with high productivity in the model tend to expect to remain only a few years in the same rental units, which means they tend to rent in the decontrolled market. Thus, they would gain from the lower free market rents following decontrol. We find that the changes in house prices and property taxes would have negligible impact on their welfare 

To sum up, full decontrol would increase average rents, but reduce initial rents of new tenants. Contrary to conventional wisdom, the effect of lower initial rent would mainly benefit newborns of low income. High-income newborns also benefit from rent decontrolling because they tend to rent in the free market.

\subsection{Benchmark vs. full rent control}\label{subsection:full_control}

In this experiment, we compare the long-run equilibrium under full rent control with the benchmark, highlighting the impact of fully removing exemptions on aggregate variables such as rents and the quantity of rental housing supplied. Because distributional outcomes are a mirror image of the previous experiment's, we leave them out.
Z's
The third column of Table \ref{tab:aggregate_effects} presents changes in aggregate variables in the equilibrium under full rent control relative to the benchmark. Without exemptions, the marginal rent would become the average rent in the controlled market, which, in the benchmark, is 10.2\% lower than the free market rent. Thus, to supply the same quantity of rental housing as in the benchmark, the initial rent of new tenants would have to rise substantially. As the initial rent of new tenants would rise to clear the rental market, the average rent would rise as well and renters' lifecycle rent profile would become more expensive. This, in turn, would increase homeownership and reduce rental demand. The decline in rental demand would outweigh the increase in homeownership, and total housing demand would decrease, thus reducing house prices. 

Despite the benefits of easier access to homeownership and lower proportional maintenance costs due to lower house prices, the higher rents would hurt the vast majority of newborns and reduce newborns' welfare. In contrast, full rent control would slightly raise average welfare, but hurt virtually all households. How is that possible? It is possible because we compute average welfare using the stationary distribution associated with the policy change in the long run, which accounts for changes in households' choices such as savings, tenure, and housing. When we compute average welfare using the benchmark stationary distribution, which does not account for changes in households' choices, average welfare actually declines, thus explaining the low approval rates. Figure \ref{fig:diagram_full_control} presents a qualitative diagram to help the reader visualize the mechanisms expounded upon in this section. 

\subsection{Optimal degree of regulation}\label{subsection:optimal_control}

Full decontrol is superior to the benchmark, whereas full control is inferior. However, an astute reader may ponder if some degree of rent control would be better. We answer this by performing a grid search to find the welfare-maximizing degree of rent control. We increase and decrease the share of the controlled market and compute the welfare associated with the long-run stationary equilibrium for each case. Figure \ref{fig:optimal_control} presents our results.

We find no degree of rent control is optimal, which means full decontrol maximizes welfare in the long run. What explains this finding? As the share of the controlled market declines, more renters are incentivized to stay less time in the same rental unit than in the benchmark. This increases homeownership rate (panel (b)) but decreases total rental demand. The latter lowers free market rents and initial rents of new tenants, and because lower initial rents benefit newborns, welfare is maximized when the market is fully decontrolled.

\subsection{Transition from benchmark to full decontrol}\label{subsection:transition}

Full decontrol is superior in the long run to all long-run equilibria in which part of the rental market is under rent control. But policymakers are often concerned about the costs incurred during the transition from one system to another. This raises the question: how long would it take for welfare gains to materialize? Answering this question is an immense computational endeavor, but it is an important practical policy question. Thus, we simulate the transition from the benchmark to the long-run equilibrium under full decontrol. We compute the transition assuming in period 1 the government unexpectedly and permanently fully decontrols rental markets.\footnote{We considered simulating a transition in which the government fully decontrols the rental market, but allows incumbent renters who were previously protected by rent control to keep their old rents until they move out, at which point the unit becomes uncontrolled. However, the computational demands of such a simulation are too immense to actually compute it.}

We present our results in panels (a)--(d) of Figure \ref{fig:transition_path}.\footnote{We do not show homeownership in transition because homeownership goes up almost immediately to its long-run value and provides no interesting insights.} In panel (a), the y-axis on the left-hand side is the scaling for the housing price, where we normalize the benchmark housing price to 1, and the y-axis on the right-hand side is the property tax rate. In panel (b), the y-axis on the left-hand side is the scaling for free market rents, where we normalize the benchmark free market rent to 1, and the y-axis on the right-hand side is the welfare gains of newborn households. In panel (c) we show the percentage of households who are better off after the policy change, and in panel (d) we present that welfare gains of households who are in retirement. In all panels, the x-axis represents periods elapsing after implementation.

Panel (b) shows free market rents would decline immediately from their benchmark level, and then decline even further as the transition continues. The decline in free market rents would be a consequence of renters reducing the size of their units. This, in turn, would cause housing prices to decline as well and necessitate a larger property tax rate to balance the budget upon shock. Why would free market rents decline? It would take households time to adjust their positions to their optimal choices under the new long-run equilibrium. Some renters would begin the period living in large units, but would no longer benefit from rent control. These renters would see an instantaneous sharp increase in their rents, which would cause them to switch to smaller units, reducing total rental housing demand, and thus decreasing free market rents beyond their long-run levels. This reduction in rents would be the primary force raising welfare immediately and all along the transition. As more households' adjust their financial position along the transition, more of them would become homeowners, and thus house prices would eventually increase slightly and the property tax rate would slightly decline. 

In panels (c) of Figure \ref{fig:transition_path}, we show that most households would benefit from the policy change. This is despite our finding that only 30\% of incumbent renters would gain upon shock, which suggests that many of the households who would gain are actually homeowners. This is because the lower free market rents would reduce the rents homeowners would have to pay upon entering a new unit should they have to go back to renting, especially during retirement when their incomes decline substantially.

In panel (d) we show that welfare gains of retirees would be negative. Why? In the benchmark, many retirees, especially those with lower incomes, have been protected by rent control for many periods. This means that even though free market rents would decline, their rents would increase relative to their benchmark protected rents, thus making them worse off.

To sum up, free market rents would decline substantially upon implementation, causing welfare gains to materialize immediately and all along the transition. Although retirees and many renters would lose, most households would gain from the lower initial rents paid upon moving into a new rental unit.

\section{Conclusion}\label{sec:conclusion}

Many cities around the world have implemented or considered implementing various types of rent controls \citep{fraser-1981}, and empirical and partial equilibrium analyses examined their impacts \citep{hayek-1929,fs-1946,fallis-1985,diamond-2019}. However, hardly any general equilibrium work explored how rent control affects the interactions between rental markets and the ownership market, and who benefits and who loses from rent control. 

We model rent control in general equilibrium with heterogeneous households. We use household-level Census data and household-level SFS data to calibrate our model to Toronto, where about 71\% of the rental market is under rent control and the rest is free. Using the calibrated model, we conducted four experiments. In the first experiment, we compared the benchmark with the long-run equilibrium under full rent decontrol. In the second, we examined the long-run equilibrium under full rent control relative to the benchmark. In the third, we searched for the welfare-maximizing degree of rent control in the long-run equilibrium. Lastly, we explored the transition dynamics from the benchmark to the long-run equilibrium under full decontrol. In all our experiments, we kept the tax code fixed, and balanced the budget using property taxes.

Our experiments revealed several key findings. Fully decontrolling rental markets would reduce initial rents paid by new tenants, but raise average rents. The former would benefit newborns, while the latter would hurt renters that have lived in controlled units for a while. Full decontrol would also substantially reduce renters' tenancy duration. The transition to full rent decontrol would feature an immediate decline in free market rents, which would raise welfare upon implementation. In contrast, full rent control would reduce the rental company's marginal rent, thus reducing the quantity of rental housing supplied and raising initial and average rents in the controlled market. As a result, welfare would decline and almost all newborns would lose. This finding suggests exemptions mitigate the negative effect of full rent control on the cost of renting and the quantity of rental housing supplied. Finally, no degree of rent control is optimal in the long run, which means full decontrol would maximize newborns' welfare in the long run. When less of the rental market is under control, more renters are encouraged to switch more quickly to homeownership. This reduces total rental demand, and thus free market rents. This, in turn, reduces initial rents of new renters and benefits newborns.

To our knowledge, our study is the first to model rent control with exemptions in general equilibrium with heterogeneous households, and thus there is much room for improvement. Due to computational limitations, we abstracted from endogenous labor supply and aggregate uncertainty, and this limited the scope of our work. We think future research may want to consider these elements once computational capacity allows. Our model abstracts from endogenous labor supply. This limited our ability to account for labor supply reactions in response to rent subsidies paid for by higher labor income taxes, which may be more effective than repealing rent control altogether, in offsetting the negative impacts rent control imposes on low-income households. Our model also abstracts from aggregate uncertainty. This prevented us from exploring the role of rent control as insurance for renters during periods in which free market rents exhibit substantial fluctuations around their long-run trends. This currently-understudied element of rent control may prove to offer substantial benefits to renters.

\bibliographystyle{aer}
\bibliography{reference_files/bib}

\vspace{1cm}

\begin{table}[H]
    \centering
    \caption{Rental market and length of stay in the same unit}
    \vspace{-0.75cm}
    \label{tab:tenure_and_length_of_stay}
    \input{tables/tenure_and_length_of_stay.tex}
    \vspace{-0.5cm}
    \begin{fignote2}{0.56\textwidth}
    \textit{Notes:} Table \ref{tab:tenure_and_length_of_stay} shows for renters in the controlled and uncontrolled market their length of stays in the same rental unit. The first row, for example, says that 17.37\% of renters in the controlled market have been in their unit less than 1 year, 33.91\% have been in their unit 1–5 years, and 48.72\% have been in their units more than 5 years. Our estimates are based on the 2016 Census for Toronto CMA.
    \end{fignote2}
\end{table}
    
\begin{table}[H]
    \centering
    \caption{Externally assigned parameter values}
    \vspace{-1.25cm}
    \label{tab:calibration_external}
    \input{tables/calibration_external}
    \vspace{-0.5cm}
    \begin{fignote2}{1.11\textwidth}
    \textit{Notes:} Table \ref{tab:calibration_external} shows the values of the externally calibrated parameters. Part (i) shows the external calibration for non-housing parameters. Part (ii) lists the external calibration of housing parameters. Part (iii) displays the external calibration of mortgage-related parameters.
    \end{fignote2}
\end{table}

\begin{table}[H]
    \centering
    \caption{Internally assigned parameter values}
    \vspace{-1.25cm}
    \label{tab:calibration_internal}
    \input{tables/calibration_internal}
    \vspace{-0.5cm}
    \begin{fignote2}{1.08\textwidth}
    \textit{Notes:} Table \ref{tab:calibration_internal} shows the values of the calibrated parameters. The model is close on all targets. Data for the percentage of homeowners with mortgages and aggregate bequests to aggregate net wealth ratio come from the 2019 SFS. Data for the total rent to total income ratio, the average income tax rate, the labor income Gini coefficient, percentage of renters in the controlled market, percentage of renters spending over 30\% of their income on rent, the homeownership rate, and the percentage of 26-34 year old households living with their parents, all come from the 2016 Census for Toronto CMA. The free market rent to housing price ratio comes from our calculations using 2016 Census and 2016 CREA data for Toronto CMA. The tax incidence of renters comes from \citet{cy-1994}, \citet{susin-2002}, \citet{fack-2006}, and \citet{sy-2020}. The homelessness rate comes from the ``Homelessness Hub''.
    \end{fignote2}
\end{table}

\begin{table}[H]
    \centering
    \caption{Aggregate non-targeted statistics}
    \vspace{-0.5cm}
    \label{tab:aggregate_non_targreted_statistics}  \input{tables/aggregate_non_targeted_statistics.tex}
    \vspace{-0.45cm}
    \begin{fignote2}{0.9\textwidth}
      \textit{Notes:} Table \ref{tab:aggregate_non_targreted_statistics} lists the performance of the model on aggregate non-targeted and distributional statistics relative to the data. The performance of the model is shown in the second column, while the third column has estimates from the literature and data. ``--'' means that the statistics is not available. 
    \end{fignote2}
\end{table}

\begin{table}[H]
    \centering
    \caption{Rental market by income quintile for all households in benchmark}
    \vspace{-0.5cm}
    \label{tab:tenure_by_income_quintile}    
    \input{tables/tenure_by_income_quintile_everyone.tex}
    \vspace{-0.45cm}
    \begin{fignote2}{0.66\textwidth}
      \textit{Notes:} Table \ref{tab:tenure_by_income_quintile} presents rental market tenure by income quintile for all households, which we compute using Toronto CMA data from the 2016 Census. Panel (a) shows statistics in the data and panel (b) the model's performance. The second column from the left represents the lowest income quintile, and the first column from the right stands for the highest quintile. The first row in each panel shows the percentage of households in each quintile who are renting in the controlled market. For example, in panel (a), 28.88\% of all households in the second quintile rent in the controlled market. The second row in each panel shows the percentage of households in each quintile who are renting in the free market. 
    \end{fignote2}
\end{table}

\begin{comment}
    
\begin{table}[H]
    \centering
    \caption{Impact of volatility, deterioration, and rent decrease on premium and welfare}
    \vspace{-0.5cm}
    \label{tab:premium_breakdown}
    \input{tables/premium_breakdown}
    \vspace{-0.45cm}
    \begin{fignote2}{0.67\textwidth}
      \textit{Notes:} Each column of Table \ref{tab:premium_breakdown} gives the impact of removing each of the components in the long-run equilibrium on the rent premium and newborns' welfare. We measure welfare in non-housing consumption-equivalent terms.
    \end{fignote2}
\end{table}

\end{comment}

\begin{table}[H]
    \centering
    \caption{Aggregate effects}
    \vspace{-0.5cm}
    \label{tab:aggregate_effects}
    \input{tables/aggregate_effects.tex}
    \vspace{-0.45cm}
    \begin{fignote2}{0.79\textwidth}
      \textit{Notes:} The second column of Table \ref{tab:aggregate_effects} compares the benchmark with the long-run equilibrium in which we impose a price ceiling on controlled units that is the same as the average controlled rent in our  benchmark and in which the proportion of the rental market that is under control is the same as in our benchmark. The third compares the benchmark with the long-run equilibrium under full decontrol. The last column compares the benchmark with the long-run equilibrium under full control. “--” means the statistic cannot be computed. Welfare is measured in non-housing consumption-equivalent terms.
    \end{fignote2}
\end{table}

\begin{comment}

\begin{table}[H]
    \centering
    \caption{Avg. years in same rental unit by labor productivity for all renters}
    \vspace{-0.5cm}
    \label{tab:avg_years_in_rental}
    \input{tables/years_in_same_rental.tex}
    \vspace{-0.45cm}
    \begin{fignote2}{1.04\textwidth}
      \textit{Notes:} Table \ref{tab:avg_years_in_rental} reports model renters' tenure length broken down by labor productivity. $e_{1}$ represents the lowest labor productivity, and $e_{7}$ the highest. In the brackets next to each productivity, we denote the percentage of households with that productivity. For example, 1.1\% of households have productivity $e_{1}$. The first row shows tenancy duration in the same unit in the benchmark where rental markets are partially regulated. The second row shows tenancy duration in the long-run equilibrium under full rent decontrol.
    \end{fignote2}
\end{table}

\end{comment}
    
\begin{table}[H]
    \centering
    \caption{Distributional effects on newborns' welfare by income quintile ($\% \Delta$ in welfare)}
    \vspace{-0.5cm}
    \label{tab:distributional_impact}
    \input{tables/distributional_impacts.tex}
    \vspace{-0.45cm}
    \begin{fignote2}{0.53\textwidth}
      \textit{Notes:} Table \ref{tab:distributional_impact} displays welfare changes for newborns, comparing price ceiling, full rent decontrol, and full control to the benchmark. Welfare is measured in non-housing consumption-equivalent terms.
    \end{fignote2}
\end{table}


\begin{figure}[H]
  \centering
    \caption{Qualitative diagram showing rent paths based on rental market choice and years in the same rental}
    \vspace{-0.37cm}
    \label{fig:diagram_rent_path}
    \includegraphics[width=\textwidth]{figures/path_of_rent_diagram.pdf}
    \vspace{-1.5cm}
    \begin{fignote2}{1.0\textwidth}
      \textit{Notes:} The downward sloping blue line shows rental prices paid by renters who live in controlled units as a function of their tenancy duration in the same units. The horizontal red line shows the average rental price faced by renters who choose to rent in the free market, and the dashed red line represents the fluctuations in free market rents. The average rental price in the free market is independent of tenancy duration. To make the intuition straightforward in this diagram, we abstract from the deterioration of units in the controlled market. Rents paid by renters in the free market are lower than rents paid by renters in the controlled market during their first years of tenancy, but higher than average rents paid in the controlled market. Why is that? If initial rents in the controlled market, $p_{r,0}$, were equal or lower than free market rents, all renters would optimally choose to rent in the controlled market. In order to shift some renters to the free market, renters in the controlled market must initially pay a premium to protect themselves against rent fluctuations and to benefit from lower rental prices over time. Renters anticipating a short stay in their rental unit and who are willing to bear the risk in the free market, would choose to rent in the free market to avoid the premium. On the other hand, renters expecting to stay many years in the same rental and whose risk tolerance is low, would pay the initial premium to enjoy the benefits of rent control.      
    \end{fignote2}
\end{figure}

\begin{figure}[H]
  \centering
    \caption{Rent-to-income ratio}
    \vspace{-0.5cm}
    \label{fig:rent_to_income_cdf}
    \includegraphics[width=\textwidth]{figures/rent_to_income_cdf.pdf}
    \vspace{-0.45cm}
    \begin{fignote2}{0.82\textwidth}
      \textit{Notes:} aaa.       
    \end{fignote2}
\end{figure}

\begin{figure}[H]
  \centering
    \caption{Lifecycle statistics}
    \vspace{-0.5cm}
    \label{fig:life_cycle}
    \includegraphics[width=\textwidth]{figures/lifecycle_plots.pdf}
    \vspace{-0.45cm}
    \begin{fignote2}{0.82\textwidth}
      \textit{Notes:} In the two panels in Figure \ref{fig:life_cycle}, we plot the performance of the model on lifecycle statistics relative to the data. In panel (a), we show the average homeownership rate for all ages. In panel (b), we show the average of the LTV ratios for all homeowners for 12 age groups because the 2019 SFS only gives the household maintainer's age in ranges of 5 years. The model output is in blue and the data is in orange. Homeownership rates by age for Toronto CMA come from the 2016 Census, and data on the average of the LTV ratios for Ontario come from the 2019 SFS.       
    \end{fignote2}
\end{figure}

\begin{figure}[H]
  \centering
    \caption{Qualitative diagrammatic representation of the impact of full rent decontrol}
    \vspace{-0.36cm}
    \label{fig:diagram_full_decontrol}
    \includegraphics[width=\textwidth]{figures/full_rent_decontrol_diagram.pdf}
    \vspace{-0.8cm}
    \begin{fignote2}{1.0\textwidth}
      \textit{Notes:} In Figure \ref{fig:diagram_full_decontrol} we provide a qualitative diagrammatic representation of the impact of full rent decontrol relative to the benchmark. Perhaps counter-intuitively, full rent decontrol would shift the rental supply curve left from the solid blue line, $\mbox{Rental supply}_{1}$, to the dashed blue line, $\mbox{Rental supply}_{2}$. Why? The increase in equilibrium house prices would increase the user-cost of housing, and thus shift the supply curve to the left. The associate decrease in the property tax would keep $\tau_{p}p$ constant, but not the interaction between house prices and the risk-free rate and the housing depreciation rate. This can be seen through equation (\ref{eq:rental_supply}), which shows that the supply of rental housing decreases when house prices increase and increases when property taxes decrease. Full rent decontrol would have two effects; one on the extensive margin (whether to own or rent), and one on the intensive margin (how much to rent for those who remain renters). Decontrolling the rental market reduces the incentive to rent and increases homeownership. This results in lower total rental demand, shifting the rental demand curve to the left, from the full red line, $\mbox{Rental demand}_{1}$, to the dashed red line, $\mbox{Rental demand}_{2}$. This shift is due to changes in the extensive margin. At the same time, decontrolling removes the rent premium for all new tenants, and this causes each remaining tenant to demand larger rental units, thus shifting rental demand from the dashed red line, $\mbox{Rental demand}_{2}$, to the dotted red line, $\mbox{Rental demand}_{3}$, which is the final location of rental demand. This second shift is due to changes in the intensive margin, which is stronger than the extensive margin. The benchmark equilibrium features quantity, $X^{1}$, average free market rent, $\overline{p}^{1}_{r,1}$, average controlled rent $\overline{p}^{1}_{r,0}$, and average rent $\overline{p}^{1}_{r}$. The equilibrium under full rent decontrol features quantity $X^{2}>X^{1}$, free market rent $\overline{p}^{2}_{r,1}>\overline{p}^{1}_{r,1}$, and average rent, which is the same as the free market rent in the new equilibrium, $\overline{p}^{2}_{r,1}>\overline{p}^{1}_{r}$.  
    \end{fignote2}
\end{figure}

\begin{figure}[H]
  \centering
    \caption{Qualitative diagrammatic representation of the impact of full rent control}
    \vspace{-0.36cm}
    \label{fig:diagram_full_control}
    \includegraphics[width=\textwidth]{figures/full_rent_control_diagram.pdf}
    \vspace{-1cm}
    \begin{fignote2}{1.0\textwidth}
      \textit{Notes:} Figure \ref{fig:diagram_full_control} is a qualitative diagram representing the impact of full rent control relative to the benchmark. Opposite to the full decontrol case, full rent control would shift the rental supply curve right from the full blue line, $\mbox{Rental supply}_{1}$, to the dashed blue line, $\mbox{Rental supply}_{2}$. Why? In contrast to the full decontrol case, here the decrease in equilibrium house prices would shift the supply curve to the right. The associate increase in the property tax would quantitatively keep $\tau_{p}p$ constant, but would not offset the interaction of house prices with the risk-free rate and the housing depreciation rate. Equation (\ref{eq:rental_supply}) demonstrates this point, showing that the supply of rental housing increases when house prices decrease and decreases when property taxes increase. Full rent control would also have two effects; the first effect would be the extensive margin effect, whereas the second effect would be the intensive margin effect. First, rental demand curve to the right from the full red line, $\mbox{Rental demand}_{1}$, to the dashed red line, $\mbox{Rental demand}_{2}$ due to changes in the extensive margin; when the whole rental market is under control, more households are encouraged to rent and stay renters longer. Second, rental demand would shift from the dash red line, $\mbox{Rental demand}_{2}$ to the dotted red line, $\mbox{Rental demand}_{3}$, which is the final location of rental demand. This shift is due to changes in the intensive margin. When the whole market is under control, all new tenants are forced to pay a premium, and thus renters who previously rented in the free market larger units, now demand smaller units. The benchmark equilibrium features quantity, $X^{1}$, free market rent, $\overline{p}^{1}_{r,1}$, average controlled rent $\overline{p}^{1}_{r,0}$, and average rent $\overline{p}^{1}_{r}$. The equilibrium under full rent control would feature quantity $X^{2}$, average controlled rent $\overline{p}^{2}_{r,0}$, and average rent, which is the same as the average controlled rent in the new equilibrium, $\overline{p}^{2}_{r,0}<\overline{p}^{1}_{r}$. In the benchmark, the quantity supplied is determined by the free market rent, $\overline{p}^{1}_{r,1}$, which is the marginal rent. Under full rent control, the quantity supplied would be determined by the average controlled rent, $\overline{p}^{2}_{r,1}$. Since $\overline{p}^{2}_{r,1}<\overline{p}^{1}_{r,1}$, the quantity supplied under full rent control, $X^{2}$, would be lower than the quantity supplied in the benchmark, $X^{1}$.
    \end{fignote2}
\end{figure}

\begin{figure}[H]
  \centering
    \caption{Optimal degree of rent control}
    \vspace{-0.5cm}
    \label{fig:optimal_control}
    \includegraphics[width=\textwidth]{figures/optimal_degree_of_control_plots.pdf}
    \vspace{-0.45cm}
    \begin{fignote2}{0.99\textwidth}
      \textit{Notes:} Figure \ref{fig:optimal_control} plots percent changes in aggregate variables as a function of the degree of rent control relative to their benchmark levels. In panel (a) we plot welfare and free market rent changes in the counterfactual relative to the benchmark. The y-axis on the left-hand side is the percent change in newborns' welfare, while the y-axis on the right-hand side is the percent change in free market rents. Welfare is measured in non-housing consumption-equivalent terms. In panel (b) plots the percentage point change in the homeownership rate. Panel (c) graphs the percent change in average rent and the quantity of rental housing supplied. The y-axis on the left-hand side is the percent change in average rents, whereas the y-axis on the right-hand side is the percent change in the quantity of rental housing supplied. In all panels, the x-axis corresponds to the degree of rent control. We show equilibrium outcomes for the degree of rent control varying from 0\% to 99\%, but not for 100\% because free market rent is not computable for that case. We discuss the 100\% case in Section \ref{subsection:full_control}.
    \end{fignote2}
\end{figure}

\newpage

\appendix

\maketitle

\counterwithin{figure}{section}
\counterwithin{table}{section}
\counterwithin{equation}{section}

\part*{Appendix (For Online Publication)}

Most of our assigned parameter values and calibration targets are taken directly from the literature. However, there are still many that we constructed ourselves from several data sources. In Section \ref{appendix:data} we describe some of the more complex ones, breaking them down to rental market variables, income-related variables, and other variables. All the variables are computed for Toronto, CMA except variables computed using the 2019 SFS, where we use data for Ontario because the lowest geography available in the publically available dataset is province.

\section{Data}\label{appendix:data}

\subsection{Rental markets}

\paragraph{Percentage of renters in same rental unit by rental market, broken down into duration bins}

The 2016 Census asks if a households rents or owns (variable name: TENUR). It also asks renters if they moved less than a year ago (variable name: Mob1) and if they moved less than 5 years ago (variable name: Mob2). The 2016 Census also asks renters when their rental unit was built (variable name: BUILTH). Thus, we can compute the fraction of renters in the controlled market who moved less than 1 year ago, 1-5 years ago, and over 5 years ago, and do the same for renters in the uncontrolled market. We show the results in Table \ref{tab:tenure_and_length_of_stay}.

\paragraph{Rent shocks in free market}

To set the rent shocks in the free market and their corresponding probabilities, we use data from the 1992-2016 Rental Market Survey (RMS).\footnote{RMS data is an internal data source not publicly available.} For each year, we compute the average rent in the controlled market and the average rent in the free market. We run the time-series for each market through an HP-filter with a penalty parameter of 1,600 to create a long-run time-series trend for each market. For each year and market we then compute the percentage deviation of the actual rent from its long-run trend. Once we have this, we compute the standard error of deviations from trend in the controlled market relative to the average rent in that market, and use it to deflate each free-market deviation. We then compute the standard error of the deflated deviations in the free market relative to the average rent in that market. Following this, we compute the number of occurrences in which the percentage deviation from trend in the free market was more than two standard errors below the trend. We then compute the average percentage deviation in those occurrences. The former allows us to compute the probability that $z$ will occur and the latter gives us the value of $z$. We repeat this for cases in which percentage deviations were 1-2, 0.5-1, and 0-0.5 standard errors below the trend, and lastly repeat the full exercise for percentage deviations above the trend.

\paragraph{Average years in same rental unit}

The 2016-2021 RMS reports the turnover rate of tenants (variable name: TURNOVER). The average turnover rate for these years is 12.25\%. With this number we use the following formula to calculate the average number of years a renter stays in the same unit:

\begin{equation}\label{eq:avg_tenure_duration}
    Y=\sum_{i=1}^{50}i(1-(1-\tau_{r})^{i}).
\end{equation}

\noindent $Y$ our notation for the average number of years renters stay in the same unit, $\tau_{r}$ is the turnover rate, and $i$ is years. To understand equation (\ref{eq:avg_tenure_duration}), consider an initial mass of renters equal to $1$. After $1$ year, a mass of $1-\tau_{r}$ still remains in their units, which means that a mass of $\tau_{r}$ lived in their unit for exactly 1 year. After $2$ years, a mass of $(1-\tau_{r})^2$ still remain in their unit, which means that a mass of $1-(1-\tau_{r})^2$ lived in their unit for exactly 2 years, and so on. We cap the calculation at 50 years because the vast majority of renters do not remain in the same unit that long. Multiplying each fraction of renters living exactly $i$ years in the same unit by the number of years they lived in their unit, gives us the expected number of years renters live in the same unit.

\paragraph{Percentage of renters in the controlled market}

To compute the percentage of renters living in controlled units, we use the 2016 Census. We use the variable BUILT to determine which units are controlled and which are not and then divided the total number of controlled units by the total number of rental units in both markets. Doing this, we find that 71\% of renters rent in the controlled market, which means that 29\% of renters rent units in the free market. We denote the fraction of renters in the controlled market with $D$.

\paragraph{Real rent decrease, $\kappa_{1}$, for each year remaining in controlled units}

In our model, we need to calibrate the rate at which real controlled rents decrease  with existing renters, which we denote in our paper by $\kappa_{1}$. To compute this, we use two data points from the Census. The first is the average rent in the controlled market (using the variables SHELCO and BUILT). The second is the average rent in the controlled of renters who moved in the past year (using the variables SHELCO, BUILT, and Mob1, and the same weights that we use to compute the average rent in the controlled market overall). We estimate that renters that moved into their controlled unit in the past year, pay $Q=13.84\%$ more than the average controlled rent in the entire market. 

Then, using the estimate of the fraction of renters renting in the controlled market, $D$, and our estimate for the average number of years renters stay in the same unit, we can estimate the average number of years a renter stays in a controlled unit, which we denote by $Y_{c}$. We denote the average number of years a renter stays in a free market unit by $Y_{f}$, and this gives us $Y=D Y_{c}+(1-D)Y_{f}$, or $Y_{c}=\frac{Y-(1-D)Y_{f}}{D}$. Given $Y_{c}$, we estimate $\kappa_{1}=0.01365$ using the following equation: 

\begin{equation}\label{eq:computing_annual_decrease}
    \kappa_{1}=((1+Q)^{\frac{1}{Y_{c}}}-1).
\end{equation}

\paragraph{Average rent to house price ratio}

To compute this statistics, we compute average annual rents across all rental units in the 2016 RMS data for, and divide it by the average apartment price in 2016 as reported in CREA data.\footnote{CREA data is an internal data source not publicly available.} This yields a ratio of $0.0524$.

\paragraph{Rent to house price ratio in free market}

To compute the rent to house price ratio in the free market, we use three numbers we estimated before. The first is our estimate of the average rent to house price ratio across both markets, which is $0.0524$ as we detailed in the previous paragraph. The second is the average share of the controlled market in the total rental market, which we estimated above to be $D=29\%$. The third is an estimate we take from \citet{zhang-gou} of how much lower the average rent in the controlled market relative to the average rent in the free market, \textsl{ceteris paribus}. 
We denote this difference with $T$. We write the average rent to price ratio as follows:

\begin{equation}\label{eq:free_market_rent_to_price}
    {\left(\frac{\overline{p}_{r}}{p}\right)}=\alpha \left(\frac{p^{0}_{r}}{p}\right)+(1-\alpha) \left(\frac{p^{1}_{r}}{p}\right).
\end{equation}

\noindent $\alpha$ is the share of the controlled market and ${\left(\frac{\overline{p}_{r}}{p}\right)}$ is the average rent to average price ratio. We know that $\frac{p^{0}_{r}}{p}=(1-T)\frac{p^{1}_{r}}{p}$, and so we can substitute this into equation (\ref{eq:free_market_rent_to_price}) and solve for $\frac{p^{1}_{r}}{p}=0.0582$.

\paragraph{Tenure distribution by income quintile}

We use 2016 Census data to compute the fraction of households in each income quintile, breaking them down into those who rent in the free market and those who rent in the controlled market. We use the total income variable (variable name: TOTINC) to compute which households occupy each income quintile. To identify renters, we use a  binary variable indicating if a household is a renter or owner (variable name: TENUR) and to identify renters in the controlled and free markets we use the categorical variable for year of dwelling construction (variable name: BUILT).

\subsection{Income and taxes}

\paragraph{Increase in real average income by age}

To compute the component of labor income that grows deterministically by age, $\zeta_{j}$, we use 2016 Census data. We compute the average market income (variable name: EMPIN) for households whose household head is between ages 25-29, and ages 55-64, respectively (variable name: AGEGRP). We then compare their incomes; households with household heads who are 55-64 years old earn 81\% more income from labor than those aged 25-29 on average.

\paragraph{Average income tax rate}

We compute the average labor income tax rate using 2016 Census data. We divided total taxes paid (variables: TOTINC-TOTINC\_AT) by total income (variable name: TOTINC). The average income tax rate we find is 19.5\%.

\paragraph{Estimating progressivity of labor income taxes, $\psi$}

In our model, we define taxes by equation (\ref{eq:taxable_labor_income}). We can rearrange this equation to obtain the following:

\begin{equation}\label{eq:taxable_labor_income_regression}
  \mbox{log}(y^{x}_{j}(e,a)-y^{\tau}_{j}(e,a))=\mbox{log}(\tau_{\ell})+\psi \mbox{log}(y^{x}_{j}(e,a)).
\end{equation}

\noindent The variable on left-hand side is after tax taxable income and the variable on the right-hand side is taxable income. The intercept is $\mbox{log}(\tau_{\ell})$ and the coefficient on taxable income is $\psi$, which determines the progressivity of taxes. We apply 2016 Census data to equation (\ref{eq:taxable_labor_income_regression}), and arrive at a point estimate for $\psi$ equal to 0.927. Since taxes in the 2016 Census are not recorded by income type, we resort to using total after tax income and total income, where our dependent variable is after tax income (variable name: TOTINC\_AT), and our independent variable is total income (variable: TOTINC).

\subsection{Other statisitcs}

\paragraph{Property tax rate}

To compute the property tax rate, we use two sources. The first is the average property tax constructed by Altus Group. The average property tax rate they find in their report is 0.66\%.\footnote{See \url{https://go.altusgroup.com/l/575253/2022-10-24/343cbn}} However, property values assessed by municipalities are significantly lower than their actual sale price. To account for this, we use house-level 1992--2016 Property Sale and Assessment Data (PSAD). In this not-publically available dataset, we can see the sale price and assessment value of each property. We use this to compute the average assessment value to sale price ratio, and find that assessment values are 33\% lower on average than sale values, implying an effective property tax rate of 0.44\%.

\paragraph{Aggregate bequests to aggregate net wealth}

To compute the aggregate bequest to aggregate net wealth ratio, we use the 2019 SFS and the Life Tables for Canada from the World Health Organization.\footnote{See \url{apps.who.int/gho/data/view.searo.60290?lang=en}.} We compute the net wealth (variable name: PWNETWPG) for economic families at each age cohort and multiply the aggregate sum by the death rate of that age cohort. This gives us the value of bequests for each age cohort. We then sum up all the bequests for all age cohorts and divide by the total sum of all economic families' net wealth. We obtain a ratio of 0.022, i.e., we estimate that aggregate bequests are 2.2\% of aggregate net wealth. Our estimate is similar to numbers estimated for the United States (see \citet{cagetti-denardi-2006}). 

\begin{comment}
    
\subsection{Other}

\paragraph{Mortgage interest rate premium}

To compute the mortgage interest rate premium, $\Delta_{m}$, we calculate the average spread between the 5-year variable rate on mortgages\footnote{This data is internal and not publically available.} and the average return on 10-year government bonds between 2002--2022 in Canada \footnote{https://www.bankofcanada.ca/rates/interest-rates/canadian-bonds/}.

\end{comment}

\end{document}
